\documentclass[11pt,a4paper]{article}
% Print layout of 002: A4 portrait, 160 x 247 mm text area.
\usepackage[a4paper,margin=25mm]{geometry}
\usepackage[T1]{fontenc}
\usepackage{xcolor}
\usepackage{lmodern,microtype,amsmath,amssymb,amsthm,booktabs,enumitem,etoolbox}
\usepackage[colorlinks=true,linkcolor=black,citecolor=magenta,urlcolor=black]{hyperref}
\urlstyle{same}
\newtheorem{thm}{Theorem}[section]
\newtheorem{prop}[thm]{Proposition}
\newtheorem{lem}[thm]{Lemma}
\newtheorem{cor}[thm]{Corollary}
\newtheorem{defn}[thm]{Definition}
\newtheorem{question}[thm]{Question}
\newtheorem{conjecture}[thm]{Conjecture}
% The four kinds of introductory statements share one counter.
\newtheorem{oldconj}{Old conjecture}
\newtheorem{introprop}[oldconj]{Proposition}
\newtheorem{introthm}[oldconj]{Theorem}
\newtheorem{newconj}[oldconj]{New conjecture}

\theoremstyle{remark}
\newtheorem{rem}[thm]{Remark}
\newtheorem{comment}[thm]{Comment}
\linespread{1.015}
\setlength{\parskip}{0.5pt}\setlength{\emergencystretch}{2em}
\setlist{itemsep=1pt,topsep=3pt,parsep=0pt}\raggedbottom
\setcounter{tocdepth}{1}
\newcommand{\topic}[1]{\par\addvspace{1.0\baselineskip}\noindent\textbf{#1.}\enspace\ignorespaces}
\title{Maximal torus theory for crossed product quantum groups}
\author{TB 2.0 and GPT 6.2}
\date{}
\makeatletter
\hypersetup{pdfauthor={\@author},pdftitle={Maximal torus theory for crossed product quantum groups}}
\makeatother
\begin{document}
\maketitle

\begin{abstract}
We give counterexamples to three proposed tests for characters,
amenability and polynomial growth using the group-dual subgroups of a
compact matrix quantum group. For an action on a compact connected Lie
group, a single such subgroup detects all three properties exactly when
the action preserves a maximal torus of the base. For a lattice
bicrossed product whose compact factor has infinitely many components,
failure of character detection is witnessed by a class function on that
factor, and a torus of a quantum quotient detects exponential growth.
These results suggest three conjectures. The character conjecture uses
the quotient by the quantum center. The amenability and growth
conjectures use tori of cocycle twists of quantum quotients.
\end{abstract}

\begingroup
\color{magenta}\hypersetup{linkcolor=magenta}
\setlength{\parskip}{0pt}
\makeatletter
\patchcmd{\l@section}{\addvspace{1.0em \@plus\p@}}{\addvspace{0pt}}{}{}
\makeatother
\tableofcontents
\endgroup

\section*{Introduction}
\phantomsection\addcontentsline{toc}{section}{Introduction}
We can read every character of a compact connected Lie group on one
maximal torus, since every conjugacy class meets it. For $G\subset U_N^+$
we use a family of tori. Throughout this paper $C(G)$ is the full
$C^*$-algebra of $G$, and $C(G)_{central}$ is the norm-closed span of
its irreducible characters. Making the fundamental matrix diagonal in
the basis given by $Q\in U_N$ defines the torus $\widehat{\Gamma}_Q$
and the restriction $\pi_Q:C(G)\to C^*(\Gamma_Q)$. We begin with
the three conjectures of \cite{bp}.

\begin{oldconj}[Characters]\label{oldchar}
If $G$ is connected, then for each nonzero $x\in C(G)_{central}$ there is a
$Q\in U_N$ such that $\pi_Q(x)\ne0$.
\end{oldconj}

For the half-classical orthogonal group $O_N^*$, restriction to the
diagonal torus distinguishes representations \cite[Theorem~6.2]{halforth}.
For half-liberations see \cite[Theorem~5.5]{bp}
and \cite[Chapter~14]{qg}.

\begin{oldconj}[Amenability]\label{oldamen}
The discrete dual $\widehat G$ is amenable exactly when all the groups
$\Gamma_Q$, $Q\in U_N$, are amenable.
\end{oldconj}

For $U_N^+$ with $N\ge2$ we need only the diagonal torus. Its group
is $F_N$, so we can already read off nonamenability. For other free examples
see \cite[Theorem~4.4]{bp} and \cite[Chapter~14]{qg}.

\begin{oldconj}[Growth]\label{oldgrowth}
Polynomial growth of $\Gamma_Q$ for every $Q\in U_N$ is equivalent to
polynomial growth of $\widehat G$.
\end{oldconj}

The dual of $O_N^*$ has polynomial growth \cite[Section~9]{halforth}.
For all three tests on the half-liberation $[H]$ of a connected
conjugation-stable $H\subset U_N$, see \cite[Theorem~5.5]{bp} and the
crossed product model in \cite{halfhopf}.

We will prove in this paper that these 3 conjectures are all wrong.
The first two counterexamples use an action of the tetrahedral group.
For growth we use the quantum automorphism group of a five-dimensional algebra.

\begin{introprop}[Characters]\label{counterchar}
There is a connected crossed product $C(SO_5)\rtimes D$ of polynomial
growth for which every torus kills the same nonzero virtual character,
the discriminant of the defining representation of $SO_5$.
\end{introprop}

We let $D=\mathbb Z^3\rtimes\mathbb Z$ act through $A_4$. This action preserves no
maximal torus of $SO_5$, so the discriminant vanishes on every torus
of the crossed product
(Theorem~\ref{sofive}).

\begin{introprop}[Amenability]\label{counteramen}
For the tetrahedral rotation action of $D=\mathbb Z^3\rtimes\mathbb Z$
on $U_2^+$, the crossed product is connected and has nonamenable dual.
Every toral group is a quotient of the virtually abelian group
$\mathbb Z\times D$.
\end{introprop}

The rotations have no common axis, so the two diagonal weights coincide
on every torus (Corollary~\ref{tetra}). If we pass to the quotient
$U_2^+$, we recover its free diagonal torus.

\begin{introprop}[Growth]\label{countergrowth}
Let $G=G^+(\mathbb C\oplus M_2)$ be the quantum automorphism group
for the canonical trace. Its dual grows exponentially, yet every toral group in every
compact matrix quantum quotient of $G$ has at most linear growth.
\end{introprop}

Here every torus lies in $SO_3$, and there is no proper nontrivial
quantum quotient (Theorem~\ref{five}). But we can twist by a cocycle
to get $S_5^+$, where a pair and a triple of points give
$\mathbb Z_2*\mathbb Z_3$.

For characters we keep the tori of $G$ itself. If they miss an element
of the character algebra, can we already see the failure on the compact
factor? We say the tori \emph{jointly detect} an algebra when their
restrictions have trivial common kernel there.

\begin{introthm}[Characters]\label{introchar}
In each of the following compact matrix cases, joint detection by the
tori of $G$ is equivalent on the three algebras $C(G)_{central}$, $C(K)_{central}$
and the pullback of $C(G/Z(G))_{central}$, where $Z(G)$ is the quantum center.
\begin{enumerate}
\item $G$ is an ordinary crossed product of a compact group $K$ by an
amenable discrete group $D$.
\item $G$ is a bicrossed product of $D=\mathbb Z^r$, $r\ge1$, and a
compact group $K$ with infinitely many connected components.
\end{enumerate}
\end{introthm}

For an ordinary action we need only ask whether each conjugacy class
of $K$ meets an invariant abelian subgroup (Theorem~\ref{classcriterion}).
We treat the lattice case in Theorem~\ref{bicrosscharacters} and pass to
$G/Z(G)$ in Corollary~\ref{classicalneutral}.

\begin{introthm}[Amenability]\label{introamen}
Let $G$ be a compact matrix bicrossed product of a discrete group $D$
and a compact group $K$. Its dual is amenable if and only if every
toral group in every unitary dual $2$-cocycle twist of a compact matrix
quantum quotient of $G$ is amenable. The torus $\widehat D\subset G$
alone detects amenability.
\end{introthm}

We find this torus simply by evaluating at the identity of $K$.
Coamenability of $G$ is equivalent to amenability of $D$
\cite[Corollary~3.7]{fmp}. Proposition~\ref{twistnecessary} explains why
all the other toral groups are amenable whenever $D$ is.

\begin{introthm}[Growth]\label{introgrowth}
Let $C(G)=C(L)\rtimes D$, with $D$ finitely generated. Take for
$L$ a compact connected Lie group, $U_2^+$ with a rotation action, or
$G^+(B)$ for the canonical trace on a finite-dimensional
$C^*$-algebra $B$. In the first and third cases any action is allowed.
Then $\widehat G$ has polynomial growth exactly when every toral group in
every unitary dual $2$-cocycle twist of a compact matrix quantum quotient
of $G$ has polynomial growth.
\end{introthm}

For a compact connected Lie base, one torus suffices for all three
tests precisely when $D$ preserves a maximal torus of the base
(Theorem~\ref{joint}). We prove both analytic tests for the
three families and their cocycle twists in Theorem~\ref{twistclass}.

We can now state our three conjectures. Theorem~\ref{introchar} gives
us a way to ask the first without choosing a crossed product presentation.

\begin{newconj}[Characters]\label{newchar}
Let $G$ be a compact matrix quantum group of Kac type with quantum
center $Z(G)$. If the tori of $G$ jointly detect the norm-closed
character algebra $C(G/Z(G))_{central}$, then they jointly detect $C(G)_{central}$.
Here characters of $G/Z(G)$ are pulled back to $G$ before restriction.
\end{newconj}

We prove this for amenable chain groups in Section~\ref{seccenter}.
In Section~\ref{sectests} we treat orthogonal easy quantum groups and
free products of free unitary quantum groups. The latter can have a
nonamenable chain group.

\begin{newconj}[Amenability]\label{newamen}
Let $G$ be a compact matrix quantum group of Kac type. Its dual is
amenable exactly when $\Lambda$ is amenable for every diagram
\[
 G\twoheadrightarrow H,\qquad H\sim K,\qquad \widehat\Lambda\subset K
\]
where $H$ is a compact matrix quantum quotient and $K$ is a unitary
dual $2$-cocycle twist of $H$.
\end{newconj}

Theorem~\ref{introamen} gives us the bicrossed-product case. For the
rotation actions on $U_2^+$, we can use the torus of the base quotient
even when every toral group of $G$ is amenable
(Section~\ref{secanalytic}).

\begin{newconj}[Growth]\label{newgrowth}
Let $G$ be a compact matrix quantum group of Kac type. Its dual has
polynomial growth exactly when $\Lambda$ has polynomial growth for
every diagram
\[
 G\twoheadrightarrow H,\qquad H\sim K,\qquad \widehat\Lambda\subset K
\]
where $H$ is a compact matrix quantum quotient and $K$ is a unitary
dual $2$-cocycle twist of $H$.
\end{newconj}

Theorem~\ref{introgrowth} covers our three families of crossed products.
For lattice bicrossed products with infinitely many components in the
compact factor, we can even exhibit a lamplighter group inside a toral
group of an untwisted quotient (Theorem~\ref{latticegrowth}).

\section{Tori and crossed products}
Let $v$ be a fundamental unitary corepresentation of a compact matrix
quantum group $G$ of Kac type \cite{wcmp}.
For a unitary matrix $Q$ of the same size as $v$, the quotient
\[
 C^*(\Gamma_Q)=C(G)\big/\big\langle (QvQ^*)_{ij}:i\ne j\big\rangle
\]
is a full group $C^*$-algebra. Its diagonal coefficients are group-like
unitaries and generate $\Gamma_Q$. We call any group-dual subgroup a
\emph{torus}. The groups $T_Q(G)=\widehat{\Gamma}_Q$ are the
\emph{standard tori}, and their indexed family is the \emph{skeleton} of
the given embedding \cite[Chapter~13]{qg}. Every torus of $G$ lies in
one of its standard tori. On the discrete side this inclusion makes the
toral group a quotient of $\Gamma_Q$. Indeed, a
unitary corepresentation of a group dual splits into weight spaces
indexed by group elements. An orthonormal basis of weight vectors
therefore diagonalizes the restriction of $v$ \cite{bbd,bp}. By Tannaka--Krein reconstruction
\cite{wtk}, the intertwiners between tensor words in $v$ and $\bar v$
determine the relations on the diagonal generators in each basis
\cite[Section~6]{bp}.

Write $C(G)_{central}$ for the norm closure in $C(G)$ of the complex span of
irreducible characters. Products correspond to tensor products and the
involution corresponds to conjugation, so this span is a unital
$*$-algebra. A virtual character is an integral linear combination of
irreducible characters. For a family of tori with restriction maps
$\pi_i$, joint detection of $C(G)_{central}$ means
\[
 C(G)_{central}\mathbin{\bigcap}\left(\bigcap_i\ker\pi_i\right)=\{0\}
\]
A single torus detects $C(G)_{central}$ if its restriction is injective, or
equivalently isometric, there. Universal norms matter in this definition.
For example the characters of $\widehat D$ are the elements of $D$,
and $C(\widehat D)_{central}=C^*(D)$ even when $D$ is nonamenable.

For a discrete quantum group $\widehat G$, amenability means that $G$ is
coamenable. For growth choose a self-conjugate generating representation
$v$ containing $1$ and use the volume
\[
 b_v(n)=\sum_{r\in v^{\otimes n}}(\dim r)^2
\]
where we count each irreducible summand of $v^{\otimes n}$ once. Tensor powers of any other
such generator give comparable lengths. We say that a torus detects an
analytic property if its discrete group has that property exactly when
$\widehat G$ does. Connectedness of $G$ means that $G$ has no nontrivial
finite quantum quotient.

Let $D$ act by automorphisms $\alpha_d$ on a compact group $K$. In the
ordinary crossed product the copy of $C(K)$ has its original coproduct,
and the canonical unitaries $s_d$, $d\in D$, satisfy
\[
 s_ds_e=s_{de},\qquad s_df s_d^*=f\circ\alpha_{d^{-1}},\qquad
 \Delta(s_d)=s_d\otimes s_d
\]
We always take the full crossed product, with its universal property for
covariant representations. We assume compact matrix type. This holds for
$K$ a compact Lie group and $D$ finitely generated, since the coefficients
of a faithful representation of $K$ and finitely many $s_d$ generate.
For a closed invariant abelian subgroup $A$ put
\[
 \Gamma_A=\widehat A\rtimes D,\qquad
 (d\cdot\lambda)(a)=\lambda(\alpha_{d^{-1}}(a))
\]
A character $\lambda$ of $A$ is a unitary function, so the covariance
relation is precisely the group relation defining $\Gamma_A$.

\begin{lem}\label{abelian}
Every torus of $G$ is contained in a torus dual to $\widehat A\rtimes D$,
where $A$ is a closed $D$-invariant abelian subgroup of $K$. Every
such $A$ gives a torus of $G$.
\end{lem}
\begin{proof}
Restrict a Hopf quotient $C(G)\to C^*(\Lambda)$ to $C(K)$.
Its image is $C(A)$ for a closed subgroup $A\subset K$, which is abelian
because the image has cocommutative coproduct. Covariance makes the
kernel of this restriction $D$-invariant. Thus the quotient factors
through $C(A)\rtimes D=C^*(\widehat A\rtimes D)$. Conversely,
equivariant restriction to $A$ gives this Hopf quotient.
\end{proof}

For a compact group $K$, $C(K)_{central}$ consists of the continuous
functions on $K$ that are constant on conjugacy classes, by Peter--Weyl. The irreducible characters of the crossed product are $s_d\chi_u$
with $d\in D$ and $u\in Irr(K)$ \cite[Theorem~6.1]{fmp}. In particular
\[
 C(G)_{central}=C^*(C(K)_{central},s_d\mid d\in D)\subset C(K)\rtimes D
\]

For $D$ trivial this gives the classical algebra of class functions. For
abelian $K$ the coefficient algebra is $C^*(\widehat K\rtimes D)$ and
$G$ is its own maximal torus. In between these cases the group unitaries
belong to $C(G)_{central}$ and can act nontrivially on the class functions. Faithfulness on $C(G)_{central}$ therefore involves the norm of sums
$\sum_d f_ds_d$ with class-function coefficients.

We shall also use compact matrix quantum quotients $G\twoheadrightarrow H$.
Write $\mathcal R_G$ for the algebra spanned by the matrix coefficients
of finite-dimensional unitary representations of $G$. Pullback gives an
inclusion $\mathcal R_H\subset\mathcal R_G$.
Restriction to a quantum subgroup gives a surjection. Thus a torus of a quotient gives
the pair of maps
\[
 \mathcal R_G\ \supset\ \mathcal R_H\
                    \longrightarrow\ \mathbb C[\Lambda]
\]
We use these maps to compare the two discrete analytic properties.
Character restrictions always use tori of $G$ itself. We shall also use
the following fact. The image of a polynomial Hopf $*$-homomorphism
into a group algebra is the span of a subgroup. Indeed, if $\sum_g c_g g$ is in the image, slicing
its coproduct at $g$ puts each $c_g g$ in the image separately. The
full completion of this subgroup algebra gives a torus of the source.

\begin{defn}\label{quotients}
The quotient test for amenability holds for $G$ if amenability of
$\widehat G$ is equivalent to amenability of every $\Lambda$ in a diagram
\[
 G\twoheadrightarrow H\supset\widehat\Lambda
\]
with $H$ a compact matrix quantum quotient and $\Lambda$ finitely
generated. The quotient test for growth uses polynomial growth on both
sides of the equivalence.
\end{defn}

\section{Characters and invariant maximal tori}\label{secweyl}

Throughout this section $K$ is a compact connected Lie group.
All its maximal tori are conjugate. A maximal torus $T$ is invariant
under $D$ when every $d\in D$ sends the subgroup $T$ onto itself.
The action on $T$ may be nontrivial. The forward construction uses
the Weyl group to recover the character norm. For the converse we
construct one virtual character that vanishes on all the tori at once.

\begin{thm}\label{weyl}
For $C(G)=C(K)\rtimes D$ with $K$ a compact connected Lie group, the
tori of $G$ detect its full character algebra if and only if $D$
preserves a maximal torus $T$ of $K$. In that case restriction to
$C(T)\rtimes D$ is isometric on the character algebra. If there is no
invariant maximal torus, a nonzero virtual character of $K$ vanishes
on every torus of $G$.
\end{thm}
\begin{proof}
Suppose first that $D$ preserves $T$. Write $W=N_K(T)/T$.
Restriction identifies $C(K)_{central}$ with $C(T)^W$ \cite{btd}. The
identification is equivariant for $D$, which normalizes the action of
$W$ on $T$.

We check faithfulness in the universal norm. Start with a faithful
representation of $C(G)$ on a Hilbert space $H$. Its class functions
and group unitaries give a covariant pair $(\pi,V)$ for $C(T)^W$ and $D$.
The Weyl average $E:C(T)\to C(T)^W$ is a $D$-equivariant conditional
expectation. Complete $C(T)\otimes H$, after dividing by null vectors,
in the following inner product, taken linear in its first variable.
\[
 \langle f\otimes\xi,h\otimes\eta\rangle
     =\langle\pi(E(\bar h f))\xi,\eta\rangle
\]
Left multiplication by $C(T)$ and the operators
$f\otimes\xi\mapsto (d\cdot f)\otimes V_d\xi$ form a covariant pair on
this completion. On $1\otimes H$ the invariant functions and the group
unitaries recover $(\pi,V)$, since for $b\in C(T)^W$ one has
$b\otimes\xi=1\otimes\pi(b)\xi$. This realizes the original norm on
$C(G)_{central}$ after restriction to $C(T)\rtimes D$.

Conversely, fix a maximal torus $T$ of $K$ and a faithful unitary
representation of $K$. Let $F$ be the set of its distinct weights on $T$.
The product
\[
 P_T=\prod_{\{\lambda,\mu\}\subset F,\ \lambda\ne\mu}
                      (\lambda-\mu)^2
\]
over unordered pairs belongs to $R(T)^W$, since $W$ permutes the weights.
It is nonzero because $R(T)$ is a Laurent polynomial ring. The
restriction isomorphism $R(K)\simeq R(T)^W$ gives a nonzero virtual
character $P$ of $K$ \cite{btd}.

If $P(t)\ne0$ at $t\in T$, the eigenspaces of $t$ in the chosen
representation are precisely the weight spaces of $T$. Any element of
$C_K(t)$ preserves these spaces, so it commutes with $T$, which acts
by scalars on each. Faithfulness and $C_K(T)=T$ give $C_K(t)=T$.
By conjugacy, $C_K(x)$ is a maximal torus for every $x\in K$ with
$P(x)\ne0$.

If a $D$-invariant abelian subgroup $A$ contains such an $x$, write
$T_x=C_K(x)$. Then $A\subset T_x$ and
$T_x\subset C_K(A)\subset C_K(x)=T_x$. Thus $C_K(A)$ is a
$D$-invariant maximal torus. When no such torus exists, $P$ vanishes
on every invariant $A$. By Lemma~\ref{abelian} it vanishes on every
torus of $G$.
\end{proof}

\section{The connected center and the three detection properties}\label{secjoint}

Suppose $C(G)=C(K)\rtimes D$ with $K$ compact matrix and $D$
finitely generated.
If $\widehat G$ has polynomial growth, then so do $\widehat K$ and $D$.
The converse depends on the action. We first give a condition under which
one can read off the growth from these two factors.

If $V$ is the coefficient space of a generating
self-conjugate corepresentation containing the trivial representation,
then $\dim V^n$ is the volume of the ball of radius $n$ in $\widehat K$.
Indeed, products of coefficients are coefficients of tensor products, and
an irreducible of dimension $m$ contributes a coefficient space of
dimension $m^2$. We can equally use any finite-dimensional, self-adjoint
algebraically generating space in $\mathcal R_K$ containing $1$ to test polynomial growth. Two such
filtrations are contained in each other after a linear change of radius
\cite{dpr}.

\begin{prop}\label{length}
Suppose some finite-dimensional, self-adjoint generating space
$V\subset\mathcal R_K$ contains $1$ and is invariant under $D$.
Then $\widehat G$ has polynomial growth if and only if both $\widehat K$
and $D$ have polynomial growth. Under this hypothesis the quotient test
for growth in Definition~\ref{quotients} passes from $K$ to $G$.
\end{prop}
\begin{proof}
Choose a finite symmetric generating set for $D$ and let $B_D(n)$ be its
word ball. Put $W=V+span\{s_d:d\in B_D(1)\}$ inside the
algebraic crossed product. Moving group elements to the right in a word
of length $n$ gives the first inclusion below. Reading a word from each
factor in succession gives the second inclusion, so that
\[
 W^n\ \subset\ \bigoplus_{d\in B_D(n)} V^n s_d
       \ \subset\ W^{2n}
\]
The middle space has dimension $|B_D(n)|\dim V^n$. This proves the
growth equivalence. If all the groups in the tests for $G$ have polynomial
growth, the quotient $G\twoheadrightarrow K$ supplies every test for
$K$ and the subgroup $\widehat D\subset G$ tests $D$ itself. The last
assertion follows.
\end{proof}

For compact connected Lie groups we can remove the invariance hypothesis
on $V$. In this case one group-dual subgroup of $G$ suffices for both
analytic properties. The length-function approach to growth of more
general bicrossed products is developed in \cite{fw}.

\begin{thm}\label{centralgrowth}
Let $K$ be a compact connected Lie group and let a finitely generated
discrete group $D$ act on $K$ by automorphisms. Put $C(G)=C(K)\rtimes D$.
Let $Z=Z(K)^0$ be the identity component of the center, and form the
discrete group $L=\widehat Z\rtimes D$ using the induced action on its character
group. Then $\widehat L\subset G$, and $\widehat G$ is amenable if and
only if $L$ is amenable. Likewise,
$\widehat G$ has polynomial growth if and only if $L$ has polynomial growth.
\end{thm}
\begin{proof}
Restriction to the characteristic subgroup $Z$ is equivariant. It gives
the Hopf quotient
\[
 C(K)\rtimes D\longrightarrow C(Z)\rtimes D=C^*(L)
\]
By \cite[Corollary~6.4]{fmp}, $G$ is coamenable precisely when $D$ is
amenable. Since $L$ is an extension of $D$ by an abelian group, the same
criterion holds for $L$. Polynomial growth of $\widehat G$ passes to
$L$ by \cite[Proposition~4.1]{dpr}.

Conversely, suppose $L$ has polynomial growth. Put $S=[K,K]$. This is a
compact connected semisimple group, and multiplication
$Z\times S\to K$ is a surjection with finite kernel \cite{btd}. It is
$D$-equivariant. Pullback on functions and the identity on $D$ give an
injection of coefficient algebras
\[
 \mathcal R_K\rtimes_{alg}D
 \ \lhook\joinrel\longrightarrow (\mathcal R_Z\otimes\mathcal R_S)\rtimes_{alg}D
 \ \simeq\ \mathcal R_S\rtimes_{alg}L
\]
Here $\rtimes_{alg}$ denotes the algebraic crossed product. Its elements
are finite sums of coefficients times group unitaries.
In the last algebra $\widehat Z$ acts trivially on $S$ and $D$ acts by
its given automorphisms. The identification sends a product of a character
of $Z$ and $s_d$ to the corresponding group element of $L$.

The outer automorphism group of $S$ is finite, as follows from its
semisimple root datum \cite{btd}. Inner automorphisms preserve the
coefficient space of each irreducible. Thus the translates under $D$ of
the coefficient space of a faithful representation of $S$ span a
finite-dimensional space. Adding its adjoint and $1$ gives an
$L$-invariant generating space. The algebra $\mathcal R_S$ has
polynomial growth \cite{dpr}. Proposition~\ref{length}, now with $S$
and $L$ as the two factors, gives polynomial growth of the algebra on
the right. Its finitely generated subalgebra on the left has polynomial
growth as well.
\end{proof}

\begin{thm}\label{joint}
Let $K$ be a compact connected Lie group, let $D$ be finitely generated
and let $C(G)=C(K)\rtimes D$. The following are equivalent.
\begin{enumerate}
\item Some torus of $G$ detects $C(G)_{central}$, amenability of $\widehat G$
and polynomial growth of $\widehat G$.
\item The family of all tori of $G$ jointly detects $C(G)_{central}$.
\item The action of $D$ preserves a maximal torus of $K$.
\end{enumerate}
For each such maximal torus $T$, the group $\Gamma_T=\widehat T\rtimes D$
gives a maximal group-dual subgroup of $G$ satisfying the first condition.
\end{thm}
\begin{proof}
A faithful restriction to one torus implies joint faithfulness.
Theorem~\ref{weyl} gives the implication from the second condition to
the third, including a single virtual character in the common kernel
when the third condition fails.

Assume that $D$ preserves $T$. Put $Z=Z(K)^0$. Since $Z$ is central and
connected, $ZT$ is a compact connected abelian group, and maximality of
$T$ gives $Z\subset T$. The restriction $\widehat T\to\widehat Z$ is
surjective by duality for compact abelian groups and is equivariant for
$D$. At the level of full coefficient algebras the maps are
\[
 C(G)\ \longrightarrow\ C^*(\Gamma_T)\
             \longrightarrow\ C^*(\Gamma_Z)
\]
Here $\Gamma_Z=\widehat Z\rtimes D$ is the group $L$ of
Theorem~\ref{centralgrowth}. The first map is faithful on $C(G)_{central}$ by
Theorem~\ref{weyl}. Amenability
and polynomial growth pass from $\widehat G$ to $\Gamma_T$ under
restriction to a quantum subgroup \cite{bp,dpr}. Each property passes
from $\Gamma_T$ to its group quotient $\Gamma_Z$, and
Theorem~\ref{centralgrowth} transfers it back to $\widehat G$. This
proves both equivalences for this same $T$.

Finally suppose a group-dual subgroup contains the torus dual to $\Gamma_T$.
By Lemma~\ref{abelian} it lies in the torus dual to $\Gamma_A$ for a closed
$D$-invariant abelian $A\subset K$. Restriction of the resulting maps
to $C(K)$ gives $T\subset A$. Since $C_K(T)=T$, we have $A=T$.
This proves maximality.
\end{proof}

For $K=SU_2$ the criterion can be read directly from the trace.
Identify the maximal tori with the unoriented axes in the Lie algebra
$Lie(SU_2)\simeq\mathbb R^3$. The inner automorphism group
$SO_3$ rotates these axes. Take $D=A_4$ acting by the rotations of a
regular tetrahedron. An invariant axis would be fixed pointwise by
each rotation of order three, yet two such rotations have different
axes. In this example $Z(K)^0=1$. The discrete group supplied by
Theorem~\ref{centralgrowth} is just $A_4$, so the dual is amenable and
has polynomial growth.

A noncentral element of $SU_2$ has distinct eigenvalues
and its centralizer is the unique maximal torus containing it. If an invariant abelian
subgroup contained such an element, its centralizer would be an invariant
maximal torus, as in Theorem~\ref{weyl}. Every invariant abelian subgroup
is therefore contained in $\{I,-I\}$. Writing $\chi$ for the defining
character, we get the explicit virtual character
\[
 P=\chi^2-4,\qquad k=\begin{pmatrix}i&0\\0&-i\end{pmatrix},\qquad P(k)=-4
\]
Both central elements have $P=0$, so every torus kills $P$.

For example take $K=\mathbb T^2$ and $D=\mathbb Z$, with the generator
acting on $\widehat K=\mathbb Z^2$ by
\[
 M=\begin{pmatrix}2&1\\1&1\end{pmatrix}
\]
Both factors have polynomial growth. Here $Z=K$ and
$L=\mathbb Z^2\rtimes_M\mathbb Z$ has exponential growth. Indeed, with
$e=(1,0)$, the $2^n$ sums $\sum_{j=0}^{n-1}\epsilon_jM^je$,
$\epsilon_j\in\{0,1\}$, are distinct. Projecting onto the expanding
eigenline gives a nonzero multiple of a sum of powers of
$(3+\sqrt5)/2>2$, in which the largest differing power decides equality.
Each sum followed by the same final power $t^n$ is a word of length at
most $2n$ using the factors $e^{\epsilon_j}t$, in multiplicative
notation. The torus in Theorem~\ref{centralgrowth} thus sees growth
created by the action.

\section{Characters for amenable actions}\label{seccharacters}

We now allow an arbitrary compact group $K$, still assuming that
$C(K)\rtimes D$ is of compact matrix type. The test can require more
than one invariant abelian subgroup. Put
\[
 S=\{k\in K : [\alpha_d(k),\alpha_e(k)]=1 \text{ for all }d,e\in D\}
\]
For $k\in S$, the closed subgroup generated by its orbit is abelian
and $D$-invariant. Conversely every point of an invariant abelian
subgroup belongs to $S$. The geometric condition below is that the conjugates of $S$ cover $K$.

\begin{thm}\label{classcriterion}
Suppose $D$ is amenable. The tori of $G$ detect every nonzero element of
$C(G)_{central}$ if and only if every conjugacy class of $K$ meets a closed
$D$-invariant abelian subgroup of $K$.
\end{thm}
\begin{proof}
The set $S$ is closed, since for each pair $d,e$ the equation
$[\alpha_d(k),\alpha_e(k)]=1$ is closed. Its conjugacy saturation
$K\cdot S=\{hkh^{-1}:h\in K,\,k\in S\}$ is compact. If a class misses
$S$, choose a continuous function on the compact Hausdorff class space
$K/Ad(K)$ which equals one on this class and vanishes on the
image of $S$. Its pullback belongs to $C(G)_{central}$ and vanishes on every
torus by Lemma~\ref{abelian}.

Conversely take $x\in C(G)_{central}$. Since $D$ is amenable the full and reduced
crossed products agree. A coefficient of a norm limit is defined by
$f_d=E_K(xs_d^*)$, where $E_K$ is the conditional expectation onto $C(K)$.
For a finite sum of characters each $f_d$ is a class function, since
automorphisms of $K$ preserve class functions. Contractivity gives the
same conclusion for $x$. The coefficients determine $x$, since their
translates give all matrix entries in the regular representation.
If $x\ne0$, choose $d$ and $k$ with $f_d(k)\ne0$. Some
conjugate of $k$ belongs to one of the subgroups $A$ in the statement.
Restriction to $C(A)\rtimes D$ has $d$-coefficient $f_d|_A$, by
continuity from finite sums, and therefore detects $x$.
\end{proof}

\begin{cor}\label{shiftcharacters}
Let $F$ be a finite group and let $D=\mathbb Z^r$, $r\ge1$, act by
translations on $K=F^D$. The tori of the crossed product detect its full
character algebra if and only if $F$ is abelian.
\end{cor}
\begin{proof}
The functions at one site and the lattice generators give compact matrix
generation. The shifts of $x\in K$ commute pairwise precisely when all
its entries commute in $F$. For abelian $F$ we can thus take $A=K$.
If $F$ is nonabelian, choose $x$ whose entries include representatives
of every conjugacy class of $F$. Conjugation in $K$ acts site by site.
If a conjugate of $x$ had commuting entries, they would generate an
abelian subgroup $J$ meeting every class of $F$. A proper subgroup
cannot do so, since its conjugates cover at most
$1+[F:J](|J|-1)<|F|$ elements. Now apply Theorem~\ref{classcriterion}.
\end{proof}

For $F=S_3$ two sites suffice. Write $s$ for the sign character and
$c$ for the character of the two-dimensional irreducible representation.
At the origin and a neighboring site $e$, take the virtual character
\[
 P(x)=(1-s(x_0))(1+s(x_e)-c(x_e))
\]
Its value is $6$ when $x_0$ is a transposition and $x_e$ is a three-cycle,
and is zero otherwise, by the character table of $S_3$. No transposition
commutes with a three-cycle, even after independently conjugating these
two entries. Thus every torus kills $P$. The same crossed product has
exponential growth detected by a torus of a quotient
(Corollary~\ref{ordinaryrow}).

There is also a one-character obstruction on $S_3\times (\mathbb Z_2)^{\mathbb Z^2}$.
Let $(m,n)$ act by conjugation by $(123)^m$ on $S_3$ and by the shift
on the second factor. The projection of any invariant abelian subgroup
to $S_3$ lies in $A_3$, since $(123)$ cycles the three subgroups generated
by transpositions. Consequently $1-s$ on the first factor vanishes on
every torus.

\section{Rotations of the free unitary group}\label{secaxis}

For a quantum base the weight spaces of a representation can play
the role of the eigenlines in the $SU_2$ example. For the fundamental
matrix of $U_2^+$ there are at most two such spaces on a group dual.
This gives an axis criterion, and retaining two different weights now
produces a nonabelian free group.

Let $D$ be a finitely generated discrete group and let
$\rho:D\to PU_2=SO_3$ be a homomorphism. Write $u$ for the fundamental
matrix of $U_2^+$, so that $C(U_2^+)$ is generated by the entries of $u$
subject to unitarity of $u$ and $\bar u$. Choose a unitary lift $U_d$ of
each $\rho(d)$. The formula $u\mapsto U_d^*uU_d$ gives an action of $D$ by
quantum group automorphisms, since the scalar ambiguity in a lift cancels.
We consider Wang's full crossed product \cite{wang}
\[
 C(G)= C(U_2^+)\rtimes D,\qquad s_d u s_d^*=U_d^*uU_d
\]
Here $s_d$ are the canonical group-like unitaries. If $S$ is a finite
generating set for $D$, a fundamental matrix for $G$ is
$v=u\oplus\bigoplus_{d\in S}s_d$. A point of $\mathbb CP^1=S^2$ is a line
in $\mathbb C^2$, and its antipode is the orthogonal line. An unoriented
axis is thus an unordered pair of orthogonal lines.

\begin{prop}\label{weights}
In any Hopf quotient onto a group algebra, the image of $u$ is either $zI_2$ or has two
distinct weights on the lines of a $\rho(D)$-invariant axis. For each
invariant axis the diagonal torus in an adapted basis is dual to
$F_2\rtimes D$, where $D$ permutes the two free generators as it permutes
the ends of the axis. If there is no invariant axis, every torus factors
through the scalar quotient $C^*(\mathbb Z\times D)$, itself a toral
quotient.
\end{prop}
\begin{proof}
Let the target be $C^*(\Lambda)$. With two distinct weights the image of
$u$ is $Pa+(I-P)b$, where $P$ is a rank-one projection and $a,b\in\Lambda$.
Conjugation by the image of $s_d$ replaces $a,b$ by their conjugates.
The crossed-product relation also gives the same matrix with $P$ replaced
by $U_d^*PU_d$. Linear independence of the group elements shows that
$U_d^*PU_d$ is either $P$ or $I-P$. The two lines therefore form an
invariant axis.

Conversely, choose an orthonormal basis on an invariant axis and leave
the $s_d$ blocks of $v$ in place. Making $u$ diagonal gives two free
unitaries $a,b$. Each $U_d$ is monomial in this basis, so covariance
requires exactly the indicated permutation of $a,b$ by $s_d$. This is the
presentation of $F_2\rtimes D$. With no invariant axis, making $u$
diagonal in any basis forces its two entries to agree. Its scalar value
$z$ then commutes with $D$, and an independent generator of $\mathbb Z$
realizes all these relations in $\mathbb Z\times D$.
\end{proof}

Set $\chi=Tr(u)$. The irreducibles of $U_2^+$ are indexed by words $w$ in
two letters corresponding to $u$ and $\bar u$. Their characters $\chi_w$
are $*$-polynomials in $\chi$, by the fusion rules \cite[Theorem~6.20]{qg}. The
irreducibles of $G$ are indexed by $(d,w)$, with character $s_d\chi_w$
\cite[Theorem~6.1]{fmp}. In particular
\[
 C(G)_{central}=C^*(\chi,s_d\mid d\in D)\ \subset C(G)
\]
The covariance relation gives $s_d\chi s_d^*=\chi$.

\begin{thm}\label{axis}
The tori of $G$ detect every nonzero element of $C(G)_{central}$ if and only if
$\rho(D)$ preserves an unoriented axis. With an axis, restriction to its
single torus is isometric on $C(G)_{central}$.
\end{thm}
\begin{proof}
If there is no axis, every torus sends $\chi$ to $2z$ by
Proposition~\ref{weights}. Hence $\chi^*\chi-4$ vanishes on all tori.
It takes the value $-4$ in the copy of $C(U_2^+)$ inside the crossed
product when we evaluate at
\[
 \begin{pmatrix}1&0\\0&-1\end{pmatrix}\in U_2
\]

Fix an invariant axis and write $\sigma(d)=0$ or $1$ according as $d$
fixes or exchanges its ends. Restriction to its torus sends
$\chi/2$ to $(a+b)/2$. To check the norm, represent $C(G)$ faithfully
on a Hilbert space $K$. Let $T$ be the image of $\chi/2$ and $V_d$ the
image of $s_d$. Then $T$ is a contraction, and $V_d$ commutes with
both $T$ and $T^*$. On $K\oplus K$ put
\[
 A=\begin{pmatrix}T&\sqrt{I-TT^*}\\
                   \sqrt{I-T^*T}&-T^*\end{pmatrix},\qquad
 J=\begin{pmatrix}I&0\\0&-I\end{pmatrix},\qquad B=JAJ
\]
The identity $T\sqrt{I-T^*T}=\sqrt{I-TT^*}T$ gives unitarity of $A$.
The assignments
\[
 a\mapsto A,\qquad b\mapsto B,\qquad
 d\mapsto (V_d\oplus V_d)J^{\sigma(d)}
\]
represent $F_2\rtimes D$. Indeed $J$ exchanges $A$ and $B$ by conjugation,
and $V_d\oplus V_d$ commutes with both. We have
\[
 (A+B)/2=T\oplus(-T^*)
\]
The first copy of $K$ therefore reduces the images of $(a+b)/2$ and
of all $d$. On this copy they act as $T$ and $V_d$. Thus the original
faithful representation of $C(G)_{central}$ is a subrepresentation of the toral
one on these generators. This proves isometry after norm closure.
\end{proof}

The analytic conclusions use two facts about the ambient $G$ which hold
for every action in this family.

\begin{prop}\label{ambient}
The quantum group $G$ is noncoamenable and $\widehat G$ has exponential
growth. If $D$ is torsion-free, then $G$ is connected.
\end{prop}
\begin{proof}
The diagonal torus of $U_2^+$ is dual to $F_2$, so $U_2^+$ is
noncoamenable \cite{bp}. The crossed-product coamenability criterion
\cite[Corollary~6.4]{fmp} gives noncoamenability of $G$.

Use word length in irreducibles with respect to $v\oplus\bar v$. The ball
of radius $n$ contains all $2^n$ distinct irreducibles $(1,w)$ with
$|w|=n$. Indeed $r_{xy}$ is a constituent of $r_x\otimes r_y$ for every
pair of words $x,y$ \cite[Theorem~6.20]{qg}. Since the volume counts each irreducible
with the square of its dimension, this gives exponential growth.

For connectedness, consider a nontrivial irreducible $(d,w)$. If $d\ne1$
and $D$ is torsion-free, its tensor powers have pairwise distinct degrees
$d^n$ in the crossed-product fusion rules \cite[Remark~6.3]{fmp}. If $d=1$,
then $w$ is nonempty and its $n$-th tensor power contains $(1,w^n)$.
In either case the coefficients generate an infinite-dimensional algebra,
which gives connectedness.
\end{proof}

\begin{cor}\label{three}
Suppose $D$ is torsion-free and has polynomial growth. If $\rho(D)$
preserves an axis, its associated torus detects $C(G)_{central}$, nonamenability
and exponential growth. In the absence of an invariant axis every
toral group has polynomial growth and every torus kills $\chi^*\chi-4$.
\end{cor}
\begin{proof}
With an axis, one toral group contains $F_2$ as the subgroup generated by
$a,b$, so it is nonamenable and has exponential growth. Theorem~\ref{axis}
gives character detection. Without an axis every toral group is a quotient
of $\mathbb Z\times D$, hence amenable and of polynomial growth.

Together with Proposition~\ref{ambient} this proves the assertions.
\end{proof}

\begin{cor}\label{tetra}
There is a connected $G\subset U_4^+$ with nonamenable dual of exponential
growth for which every toral group has polynomial growth and every torus
kills the same nonzero virtual character.
\end{cor}
\begin{proof}
Take $D=\mathbb Z^3\rtimes\mathbb Z$, where the generator $t$ cycles the
standard basis of $\mathbb Z^3$. This group is torsion-free, and
$\mathbb Z^3\times\langle t^3\rangle$ has index three. Map the basis to
the three nonidentity elements of the Klein four group in $A_4$, in the
order in which conjugation by $(123)$ cycles them, and send $t$ to $(123)$.
This gives a surjection onto the tetrahedral rotation group $A_4\subset SO_3$.
An invariant pair of antipodal points would be fixed pointwise, since
$A_4$ is generated by elements of order three. Two tetrahedral rotations
with distinct axes have no common fixed point. Corollary~\ref{three} now
applies. The elements $e_1,t$ generate $D$, so the fundamental matrix
$u\oplus s_{e_1}\oplus s_t$ has size four.
\end{proof}

For compact Lie groups and for group duals, the skeleton determines the
embedded group, and every proper subgroup in the same class has a strictly
smaller standard torus for some $Q$ \cite[Theorem~13.12]{qg}.
The tetrahedral action lets us compare two skeletons in the same $U_4^+$.

\begin{thm}\label{skeleton}
Let $G\subset U_4^+$ be the tetrahedral example of Corollary~\ref{tetra}
and let $H$ be its scalar torus, dual to $\mathbb Z\times D$. Then
$H\subsetneq G$ and, as embedded subgroups of $U_4^+$,
\[
 T_Q(H)=T_Q(G)\qquad(Q\in U_4)
\]
The discrete group $\widehat H$ has polynomial growth, whereas
$\widehat G$ is nonamenable and has exponential growth.
\end{thm}
\begin{proof}
Equip $H$ with the restriction of the fundamental matrix of $G$.
For any inclusion $H\subset G$ with this convention, forming the
diagonal quotient in the basis $Q$ imposes the same off-diagonal
relations on both sides. Thus $T_Q(H)=H\cap T_Q(G)$. In our example
Proposition~\ref{weights} gives $T_Q(G)\subset H$ for every $Q$, which
proves the asserted equality. In the original block basis $T_1(G)=H$.
The element $\chi^*\chi-4$ from
the proof of Theorem~\ref{axis} is nonzero on $G$ and vanishes on $H$,
proving properness. Finally $\mathbb Z\times D$ is virtually abelian,
and Proposition~\ref{ambient} gives the assertions about $\widehat G$.
\end{proof}

The same tetrahedral group gives a classical base of polynomial growth
with an explicit character in the common kernel of all restrictions.
\begin{thm}\label{sofive}
There is a connected compact matrix quantum group of Kac type, generated
by two group-dual subgroups, for which all tori kill the same nonzero
virtual character. Its discrete dual has polynomial growth.
\end{thm}
\begin{proof}
Let $K=SO_5$. Let $D=\mathbb Z^3\rtimes\mathbb Z$ act through the
surjection $D\to A_4$ of Corollary~\ref{tetra}. Embed $A_4$ in $SO_5$
by even permutations of the first four coordinates and use conjugation.
Form $C(G)=C(K)\rtimes D$. This $G$ is connected. Indeed, an
irreducible of nonidentity $D$-degree has tensor powers of infinitely
many degrees, and in identity degree its coefficients belong to $C(K)$,
where a nontrivial irreducible has infinite image by connectedness.
The coefficients of the defining representation of $SO_5$, together
with $1$, span a $D$-invariant generating space $V$. Both the dual of $SO_5$ and the virtually
abelian group $D$ have polynomial growth \cite{dpr}.
Proposition~\ref{length} applied to $V$ gives the growth assertion.

A maximal torus of $SO_5$ rotates two orthogonal real planes and fixes
their orthogonal line. Its normalizer preserves this decomposition up to
interchanging the two planes. If $A_4$ normalized such a torus, it would
preserve each plane, since it has no quotient of order two. The same
fact makes its determinant on each plane positive and its action on
the line trivial. Its image would lie in $SO_2\times SO_2$, contrary
to the faithful nonabelian permutation action.

The defining representation of $SO_5$ has five distinct weights, each
of multiplicity one. The virtual character in Theorem~\ref{weyl} is
therefore the usual discriminant
\[
 P(k)=disc_X\det(XI_5-k)
\]
For instance it is nonzero at a rotation with eigenvalues
$1,e^{\pm\pi i/3},e^{\pm\pi i/2}$. Its expression as a polynomial in
the characteristic coefficients writes it directly in terms of exterior
power characters. By Theorem~\ref{weyl} it vanishes on every torus of $G$.

It remains to show generation. Let $E$ be the diagonal sign subgroup
of $SO_5$, and let $C$ be the circle of rotations in the plane spanned
by $e_5$ and $v=(e_1+e_2+e_3+e_4)/2$. The group $A_4$ preserves $E$
and fixes $C$ pointwise. The conjugates of the infinitesimal rotation
$e_5\wedge v$ by $E$ span all $e_5\wedge e_i$, $1\le i\le4$. Indeed
the first four signs are arbitrary and the fifth contributes an overall
sign. Bracketing these rotations gives the remaining $e_i\wedge e_j$.
It follows that finite products from $E$ and $C$ are dense in $SO_5$.

Take the tori with discrete groups $\widehat E\rtimes D$ and
$\widehat C\rtimes D$ from Lemma~\ref{abelian}. Suppose a Hopf
$*$-ideal lies in the kernels of both restrictions and take an element
$\sum_d f_ds_d$ of this ideal in $\mathcal R_G$. Its iterated
coproduct vanishes after restriction to any ordered list of the two
tori. The coefficient of $(d,\ldots,d)$ in that restriction is
$(a_1,\ldots,a_m)\mapsto f_d(a_1\cdots a_m)$, with $a_i$ in the
corresponding $E$ or $C$. Density forces each $f_d$ to vanish. Hence
these two tori generate $G$, even though the discriminant vanishes on
each of them separately.
\end{proof}

\section{Tori of quotients and quantum automorphism groups}\label{secquotients}

In the rotation examples of Section~\ref{secaxis}, pullback of the
coefficients of $U_2^+$ gives a quantum quotient
$G\twoheadrightarrow U_2^+$. Its diagonal torus is dual to $F_2$ for
every action of $D$. The following observation explains when
the quotient test passes from a base to its crossed product.
\begin{prop}\label{permanence}
Suppose the quotient test for amenability in Definition~\ref{quotients} holds
for a compact matrix quantum group $K$ of Kac type. It also holds for
$G$ with $C(G)=C(K)\rtimes D$, for any action of a finitely
generated discrete group $D$ by quantum group automorphisms.
\end{prop}
\begin{proof}
Suppose every group tested for $G$ is amenable. The inclusion
$C(K)\subset C(G)$ gives a quantum quotient $G\twoheadrightarrow K$.
Composing with this quotient shows that every group tested for $K$ is
also tested for $G$, so $K$ is coamenable by hypothesis.

There is a second test inside $G$ itself. The counit of $K$ is invariant
under the action of $D$, so evaluation at the counit on $C(K)$ and
$s_d\mapsto d$ give a Hopf $*$-surjection $C(G)\to C^*(D)$.
Thus $\widehat D\subset G$ and $D$ is amenable. By the crossed-product
criterion \cite[Corollary~6.4]{fmp}, $G$ is coamenable. In the
other direction, coamenability passes to compact quantum quotients
and then to their tori, as recalled in Proposition~\ref{twistnecessary}.
\end{proof}

For polynomial growth Proposition~\ref{length} gives an analogous
conclusion when the action preserves a finite-dimensional generating
space. We next examine how much these quotient tests can see for quantum
automorphism groups. The exceptional algebra occurs already in dimension
five.
Consider the quantum automorphism group $G$ of the algebra
$B=\mathbb C\oplus M_2(\mathbb C)$ with its canonical trace
\[
 \tau(z,X)=\frac{z+2Tr(X)}5
\]
where $Tr$ is the ordinary matrix trace. This is the normalized trace of
the left regular representation. The universal trace-preserving coaction
on $B$ gives a compact matrix quantum group of Kac type \cite{generic}.
Conjugating the matrix summand gives a classical subgroup
$PU_2=SO_3\subset G$.

\begin{lem}\label{fiveorbit}
Every group-dual subgroup of $G$ is contained in this $SO_3$.
The possible discrete groups are precisely $\mathbb Z$, the finite
cyclic groups, and the Klein four group.
\end{lem}
\begin{proof}
A group-dual subgroup gives a grading $B=\bigoplus_g B_g$ and a coaction
$\alpha(x)=x\otimes g$ for $x\in B_g$. Let $p:B\to\mathbb C$ be the
scalar character. The algebra map
\[
 \varphi=(p\otimes id)\alpha:B\longrightarrow\mathbb C[\Gamma]
\]
sends $x\in B_g$ to $p(x)g$. Its image is spanned by the group elements
which occur. Since it is a finite-dimensional unital $*$-algebra, these
elements form a finite subgroup $F\subset\Gamma$ and the image is
$\mathbb C[F]$. If $\varphi$ were injective, $|F|=5$. Such a group is
cyclic, whereas $B$ is noncommutative. Hence $\ker\varphi\ne0$.
Composing $\varphi$ with the counit gives $p$, so
$\ker\varphi\subset\ker p=M_2(\mathbb C)$. Simplicity of the matrix
summand gives equality.

The map $\varphi$ preserves degrees, so the matrix summand is a graded
ideal. If $e$ is its unit, comparison of homogeneous components in
$ex=x$ shows that the neutral component $e_1$ acts as the identity
on every homogeneous element of the ideal. Thus $e_1e=e$, while
$e_1e=e_1$ since $e_1$ belongs to the ideal. Both $e$ and $1-e$ are
therefore fixed. On
$M_2(\mathbb C)$ the coaction preserves the ordinary trace. The universal
trace-preserving quantum automorphism group of $M_2(\mathbb C)$ is
$PU_2=SO_3$. Indeed the quantum and classical groups here have the
$SO_3$ fusion rules \cite{generic}, and restriction sends the
three-dimensional summand of $M_2$ to the usual rotation representation.
Induction with the fusion rules identifies every irreducible on
restriction, hence also its coefficient coalgebra. This proves the
asserted factorization.

Since $C(SO_3)$ is commutative, every group-dual subgroup is classical
abelian. A rotation through an angle other than $0$ or $\pi$ has a circle
as its centralizer. A closed
subgroup of this circle is finite cyclic or the full circle. If every
nonidentity element is a half-turn, commuting half-turns about distinct
axes have perpendicular axes, and two of them generate a Klein four
group. This gives the list on taking duals, and each case occurs.
\end{proof}

\begin{thm}\label{five}
The compact quantum group $G$ has no proper nontrivial compact quantum
quotient. Its discrete dual is nonamenable and has exponential growth.
Nevertheless in every diagram $G\twoheadrightarrow H\supset\widehat\Lambda$
the group $\Lambda$ has at most linear growth. In particular the two
quotient tests in Definition~\ref{quotients} both fail, even for connected $G$.
\end{thm}
\begin{proof}
The irreducibles of $G$ have the $SO_3$ fusion rules \cite{generic}.
Writing them as $r_0=1,r_1,r_2,\ldots$, the coaction on $B$ is $1+r_1$ and
\[
 r_j\otimes r_k=\bigoplus_{s=|j-k|}^{j+k}r_s
\]
Under a compact quantum quotient the polynomial Hopf algebra of $H$
embeds in that of $G$. Its simple coefficient coalgebras give a subset
of these irreducibles closed under tensor products and constituents.
If it contains $r_j$ with $j\ge1$, then $r_j\otimes r_j$ contains $r_1$,
and tensor powers of $r_1$ contain every $r_k$. The inclusion is thus
either trivial or onto, as in Wang's proof in \cite{simple}.
In particular $G$ has no nontrivial finite quantum quotient, so is connected.

For growth, put $d_k=\dim r_k$. Since $\dim B=5$, the fusion rule gives
\[
 d_0=1,\qquad d_1=4,\qquad d_{k+1}=3d_k-d_{k-1}\quad(k\ge1)
\]
In particular $d_{k+1}\ge2d_k$ by induction. The ball of radius $n$ for
$1+r_1$ has volume $\sum_{k=0}^n d_k^2$, which grows exponentially.
Nonamenability in dimension five follows from the coamenability
criterion of \cite{generic}. Finally $H$ is trivial or isomorphic to $G$,
so Lemma~\ref{fiveorbit} accounts for every $\Lambda$ in the test.
\end{proof}

This is a nonclassical quantum group all of whose group-dual subgroups
are classical. Such a group is called \emph{strange} in \cite{bbd}. The same pair of matrix sizes also occurs in the computation of
outer quantum symmetries in \cite[Section~3.3, Example~2]{goswami}.

\begin{cor}\label{fivecharacters}
Let $G=G^+(\mathbb C\oplus M_2)$ with its canonical
trace. Restriction to one circle is injective on the algebraic span of
the irreducible characters, whereas every torus kills the same nonzero
element of the norm-closed character algebra $C(G)_{central}$.
\end{cor}
\begin{proof}
Put $x=\chi_{r_1}=x^*$. The fusion rules in Theorem~\ref{five} express
every irreducible character as a polynomial in $x$. On the subgroup
$SO_3$ the action on $\mathbb C\oplus M_2$ splits as two trivial
representations and the three-dimensional rotation representation.
Hence $x|_{SO_3}=1+Tr(R)$. On rotations through angle $\theta$
about a fixed axis this is $2+2\cos\theta$, with range $[0,4]$.
Thus a polynomial in $x$ vanishing on this circle is the zero polynomial.

Let $h$ be the Haar state. Irreducibility and the fusion rule give
$h(x)=0$ and $h(x^2)=1$. If $x$ were positive, the inequality
$x^2\le\|x\|x$ would contradict these two equalities. Therefore
$x_-=\max(-x,0)$ is a nonzero element of $C(G)_{central}$. Its restriction to
$SO_3$ is zero, and all tori factor through that subgroup by
Lemma~\ref{fiveorbit}.
\end{proof}

Is this failure specific to the algebra of dimension five? For the
canonical trace, the possible block sizes give a complete answer.

\begin{prop}\label{onlyfive}
Among the quantum groups $K=G^+(B)$ with the canonical
trace, $B=\mathbb C\oplus M_2$ is the unique algebra, up to isomorphism,
for which the tori of $K$ fail to detect amenability or polynomial growth.
\end{prop}
\begin{proof}
Put $N=\dim_{\mathbb C} B$. For $N\le3$ we have the finite groups $S_N$. At $N=4$ the irreducibles
have the $SO_3$ fusion rules with $\dim r_1=3$ and hence
$\dim r_j=2j+1$. The ball volumes $\sum_{j=0}^n(2j+1)^2$ are polynomial.
Both cases are coamenable \cite{generic}, and both properties pass to the
tori \cite{bp,dpr}.

Let $N\ge5$. We produce a trace-preserving group grading whose support
generates a nonamenable group of exponential growth. Its coaction on $B$
is faithful in the sense that its coefficients generate the group algebra,
so its group dual embeds in $K$. By permanence this forces the ambient
dual to be nonamenable and to have exponential growth. For a summand $M_k$ one can give matrix
units $e_{ij}$ degree $g_i^{-1}g_j$ for any list of group elements $g_i$.
On a scalar block $\mathbb C^m=C^*(\mathbb Z_m)$ one can use the Fourier
grading by $\mathbb Z_m$. Both gradings preserve the canonical trace.
Gradings on distinct direct summands combine as gradings by the free
product of their groups, since products between summands vanish.
All unused summands may be placed in the identity degree.

If $B$ has a summand $M_k$ with $k\ge3$, use the list $1,a,b,1,\ldots,1$
in $F_2=\langle a,b\rangle$. If it has two $M_2$ summands, grade each
by $\mathbb Z$, giving $F_2$ once again. If there is just one $M_2$ and
all other summands are scalar, two scalar summands together give
$\mathbb Z_2$, and the matrix summand gives $\mathbb Z$. This gives
$\mathbb Z_2*\mathbb Z$ whenever there are at least two scalar summands.
Finally, for $B=\mathbb C^N$, splitting off a pair and a triple gives
$\mathbb Z_2*\mathbb Z_3$. Each displayed group contains a free group on
two generators. The sole unaccounted case is $\mathbb C\oplus M_2$,
which fails both tests by Theorem~\ref{five}.
\end{proof}

The exceptional algebra has nonamenable dual. Its failure of the quotient
growth test leaves the following implication to be determined.

\begin{question}\label{amenablegrowth}
Suppose $\widehat G$ is amenable and every group $\Lambda$ in a diagram
$G\twoheadrightarrow H\supset\widehat\Lambda$ as in
Definition~\ref{quotients} has polynomial growth. Must $\widehat G$ have
polynomial growth?
\end{question}

\section{Two conjectures on amenability and growth}\label{secanalytic}

For $G=G^+(\mathbb C\oplus M_2)$ there is another way
to recover a torus of exponential growth. A cocycle representing the nontrivial
class on the Klein four group twists $G$ into $S_5^+$ \cite[Theorem~4.8]{mrozinski}.
The latter contains the dual of $\mathbb Z_2*\mathbb Z_3$, acting on a pair
and a triple of points \cite{bbd}. More generally, for a finite-dimensional
$C^*$-algebra $B$ with its canonical trace, $G^+(B)$ is a
unitary cocycle deformation of $S_N^+$, where $N=\dim_{\mathbb C}B$
\cite[Theorem~3.1]{behy}.

A unitary dual $2$-cocycle twist gives a unitary monoidal equivalence
which preserves the Hilbert space dimension of every representation. Conversely, a dimension-preserving unitary
monoidal equivalence is realized, up to isomorphism, by such a twist
\cite{bdv}. Write $H\sim K$ for this equivalence. The identity quotient and the trivial twist are allowed throughout.

\begin{prop}\label{twistnecessary}
Let $G$ be a compact matrix quantum group of Kac type. For every diagram
\[
 G\twoheadrightarrow H,\qquad H\sim K,\qquad \widehat\Lambda\subset K
\]
with $H$ a compact matrix quantum quotient of $G$ and $K$ a unitary dual
$2$-cocycle twist of $H$, amenability of $\widehat G$ implies amenability
of $\Lambda$, and polynomial growth of $\widehat G$ implies polynomial
growth of $\Lambda$. A dimension-preserving unitary monoidal
equivalence of $G$ preserves the collection of groups $\Lambda$ arising
in these diagrams, up to isomorphism.
\end{prop}

\begin{proof}
Growth passes through compact quantum quotients and subgroups
\cite{dpr}. Under $H\sim K$ a self-conjugate generating representation
containing $1$ and its image have exactly the same constituents in
each tensor power, with the same dimensions. The two ball volumes agree
at every radius.

For a compact quantum quotient the counit restricts, and the inclusion
of coefficient algebras is isometric for reduced norms because the Haar states
agree and the reduced Haar state is faithful. Thus the quotient inherits
a bounded reduced counit. For a compact quantum subgroup, take a self-conjugate generator
$v$ containing $1$ and put $m=\dim v$. Restriction can only increase
the multiplicity of $1$ in $v^{\otimes 2n}$. This multiplicity is the
$2n$-th Haar moment of $\chi_v$, whose $2n$-th root tends to the
reduced norm by faithfulness of the reduced Haar state. Hence if the
ambient norm is $m$, the subgroup norm is also $m$. On both sides we can write $v=1\oplus w$, with
$\chi_v=1+\chi_w$ and $\|\chi_w\|\le m-1$. Thus the spectrum is
contained in $[2-m,m]$. Norm $m$ is therefore equivalent to the
presence of the positive endpoint $m$, the quantum Kesten criterion
for coamenability \cite{bp}.
Under $H\sim K$ both $m$ and the multiplicity of $1$ in every
$v^{\otimes 2n}$ are preserved. Thus the reduced character norm attains
$m$ on one side exactly when it does on the other. This also follows
from the fusion-ring criterion \cite[Theorem~4.5]{kyed}.

Finally a compact matrix quantum quotient is specified by a full rigid
tensor subcategory of $Rep(G)$ closed under subobjects and
finite sums and generated by one object and its conjugate. Its coefficient
algebra is the span of the coefficient spaces of the simple objects in
the subcategory. A monoidal equivalence carries each such subcategory
to another and restricts to an equivalence there, preserving dimensions.
Composing the inverse of the transported subcategory equivalence with $H\sim K$ gives the same
tested group. The inverse equivalence proves equality of the collections.
\end{proof}

In particular, any finite succession of compact matrix quantum quotients
and unitary dual cocycle twists before taking a torus reduces to a
single quotient followed by a single twist. Indeed, successive quotients compose,
and a quotient taken after an equivalence can be transported back along
that equivalence by the preceding proof.

\begin{conjecture}[Amenability]\label{twistamen}
Let $G$ be a compact matrix quantum group of Kac type. Its discrete
dual is amenable if and only if $\Lambda$ is amenable in every diagram
\[
 G\twoheadrightarrow H,\qquad H\sim K,\qquad \widehat\Lambda\subset K
\]
where $H$ is a compact matrix quantum quotient of $G$ and $K$ is a
unitary dual $2$-cocycle twist of $H$.
\end{conjecture}

\begin{conjecture}[Polynomial growth]\label{twistgrowth}
Let $G$ be a compact matrix quantum group of Kac type. Its discrete
dual has polynomial growth if and only if $\Lambda$ has polynomial
growth in every diagram
\[
 G\twoheadrightarrow H,\qquad H\sim K,\qquad \widehat\Lambda\subset K
\]
where $H$ is a compact matrix quantum quotient of $G$ and $K$ is a
unitary dual $2$-cocycle twist of $H$.
\end{conjecture}

The open direction in each case is the existence of a witness. In the
first case a nonamenable dual should yield a nonamenable $\Lambda$. In
the second case the conclusion sought is a single $\Lambda$ with growth
that is not polynomial, even when $\widehat G$ is amenable. The group
$\mathbb Z_p\wr\mathbb Z$ supplies such a witness in Corollary~\ref{ordinaryrow}.

The permanence arguments of Propositions~\ref{permanence} and
\ref{length} also apply to Conjectures~\ref{twistamen} and~\ref{twistgrowth}. Indeed, for
$C(G)=C(L)\rtimes D$, every diagram for $L$ is a diagram for $G$
after composition with $G\twoheadrightarrow L$, and $\widehat D\subset G$
supplies the other factor. Thus the amenability implication passes from
$L$ to $G$ for any action. The growth implication does so whenever $D$
preserves a finite-dimensional generating space as in Proposition~\ref{length}.

\begin{thm}\label{twistclass}
Conjectures~\ref{twistamen} and~\ref{twistgrowth} hold for
$C(G)=C(L)\rtimes D$, where $D$ is finitely generated and acts by
quantum group automorphisms, for the following bases and actions.
\begin{enumerate}
\item $L$ is a compact connected Lie group.
\item $L=U_2^+$ with the rotation action of Section~\ref{secaxis}.
\item $L=G^+(B)$ for a finite-dimensional $C^*$-algebra
with its canonical trace.
\end{enumerate}
They also hold for every unitary dual cocycle twist of these crossed products.
\end{thm}
\begin{proof}
In the first case Theorem~\ref{centralgrowth} supplies a single group
inside $G$ testing both properties. In the second case the quotient
$G\twoheadrightarrow U_2^+$ supplies $F_2$ without twisting.

In the third case $G\twoheadrightarrow L$. If $N=\dim B\ge5$,
Proposition~\ref{onlyfive} gives a group detecting both failures directly,
except at $B=\mathbb C\oplus M_2$. In that case take
the twist $K=S_5^+$ of $L$ and its toral group $\Lambda=\mathbb Z_2*\mathbb Z_3$. Equivalently, for every $N\ge5$ one can use the
same pair and triple in the model $S_N^+$ of \cite[Theorem~3.1]{behy}.

If $N\le4$, the base $L$ is coamenable and has polynomial growth, as in
Proposition~\ref{onlyfive}. The counit supplies $\widehat D\subset G$,
so the test includes $D$. When $D$ is amenable, $G$ is coamenable by
\cite[Corollary~6.4]{fmp}. For growth, Proposition~\ref{length} applies
to any action. For $N\le3$ use the whole finite-dimensional algebra, and
for $N=4$ use the coefficient space of $1+r_1$. Indeed $r_1$ is the
unique irreducible whose tensor square has three distinct constituents,
so every automorphism preserves this space. Thus polynomial growth of
$D$ implies polynomial growth of $\widehat G$. The last assertion follows
from Proposition~\ref{twistnecessary} and invariance of the two properties.
\end{proof}

The third case also includes every compact quantum group $L$ of Kac
type having the $SO_3$ fusion rules. By \cite{mrozinski} it is
$G^+(B,\varphi)$ for a homogeneous faithful state, with
$B=1+r_1$ as a representation. On this finite-dimensional algebra the invariant state is the
canonical trace. To see this, let $U$ implement the coaction on
$L^2(B,\varphi)$. Left multiplication $b\mapsto L_b$ intertwines the
coaction with conjugation by $U$ on $End(B)$. In Kac type
$\bar U$ is unitary, so conjugation preserves the ordinary matrix trace.
Thus $b\mapsto Tr(L_b)/\dim B$ is invariant. Uniqueness follows from the
single copy of $1$ in $B$, giving the assertion.

\section{Bicrossed products with a compact Lie factor}\label{secmatched}

Let $D$ be discrete and $K$ compact. The following formulas also apply
to the infinite component groups of the next two sections. For a matched pair we use the
factorization $dk=\alpha_d(k)\beta_k(d)$, with $\alpha_d(k)\in K$ and
$\beta_k(d)\in D$. Factoring $dhk$ and $dsk$ in two ways gives
\[
 \alpha_d(hk)=\alpha_d(h)\alpha_{\beta_h(d)}(k),\qquad
 \beta_k(ds)=\beta_{\alpha_s(k)}(d)\beta_k(s)
\]
Here $\alpha$ is a left action and $\beta$ a right action. When
$\beta$ is trivial, $\alpha$ acts by group automorphisms and we
recover the ordinary crossed product. Write $G$ for
the compact bicrossed product. Its coefficient algebra is the algebraic
crossed product $\mathcal R_K\rtimes_{alg}D$, with
$s_d f s_d^*=f\circ\alpha_{d^{-1}}$. Functions have their usual
coproduct, and the canonical unitaries satisfy
\[
 \Delta(s_d)=\sum_{r\in\beta_K(d)}
 s_d 1_{\{k:\beta_k(d)=r\}}\otimes s_r
\]
The orbits are finite by continuity and compactness. The Haar state
integrates the identity coefficient against Haar measure on $K$. It
is a trace since $\alpha$ preserves this measure. Evaluation at the
identity of $K$ gives $\widehat D\subset G$, and $\widehat G$ is
amenable exactly when $D$ is amenable \cite[Theorem~3.4 and
Corollary~3.7]{fmp}. Thus for bicrossed products the remaining analytic
question concerns growth when $D$ itself has polynomial growth.

\begin{thm}\label{finitecomponents}
Let $K$ be a compact Lie group, let $D$ be finitely generated and put
$K_0=K^0$ and $Z=Z(K_0)^0$. In the bicrossed product the torus dual to
$\widehat Z\rtimes D$ detects both amenability and polynomial
growth. The action of $D$ on the character group here comes from its
action on $K_0$ by automorphisms.
\end{thm}
\begin{proof}
The continuous map $h\mapsto\beta_h(d)$ is constant on $K_0$, with
value $d$. Each $\alpha_d$ fixes the identity and preserves $K_0$.
By the matched-pair identities, its restriction $\theta_d$ to $K_0$ is an
automorphism and
\[
 \alpha_d(hk)=\theta_d(h)\alpha_d(k)\qquad(h\in K_0)
\]
Restriction of functions to $K_0$ and $s_d\mapsto s_d$ give the ordinary
crossed product on $K_0$ as a quantum subgroup. Theorem~\ref{centralgrowth}
supplies the stated torus there, and the amenability assertion follows.

For growth, choose representatives $k_1=1,\ldots,k_b$ for the cosets
$K_0k_i$. Take a finite-dimensional self-adjoint generating space
$V\subset\mathcal R_{K_0}$ containing $1$ and stable under left
and right translations. The coefficient spaces of a faithful
representation, its conjugate and $1$ give such a $V$. Use independent
copies of $V$ on the $b$ cosets to obtain a space $E$ on $K$. A
coefficient on one coset extended by zero has finite-dimensional right
translation span, since right translation permutes the cosets and acts
on their coordinates by translations. Hence $E\subset\mathcal R_K$.
Conversely, each coefficient on $K$ restricts on each coset to a
representative function on $K_0$. Since $E$ includes the coset
indicators, it generates $\mathcal R_K$.

Adjoin the unitaries for the same symmetric set of generators of $D$ to
$E$ and to $V$, giving spaces $W$ and $W_0$ in the two crossed products.
For a word of degree $d$ with coefficient $f$, record the $b$ entries
$f(hk_i)s_d$. Multiplying on the right by $s_e$ replaces $d$ by $de$.
Multiplying on the right by $v\in E$ multiplies the $i$th entry on the
right by the function $h\mapsto v(h\alpha_{d^{-1}}(k_i))$ in $V$.
Indeed, commuting this function past $s_d$ introduces $\theta_{d^{-1}}$
on $h$, giving precisely $v(\alpha_{d^{-1}}(hk_i))$. Thus the entries
of a word of length $n$ lie in $W_0^n$. Recording the entries degree by degree defines an injective linear map
into $b$ copies of $W_0^n$, since the cosets cover $K$. Restriction to the
first coset maps $W^n$ onto $W_0^n$. We obtain
\[
 \dim W_0^n\ \le\ \dim W^n\ \le\ b\dim W_0^n
\]
The polynomial growth assertion now follows from Theorem~\ref{centralgrowth}.
\end{proof}

\section{The component group and the cyclic case}\label{seccomponent}

Let $D=\mathbb Z^r$, $r\ge1$, and assume that $K/K^0$ is infinite.
Every bicrossed product in this class is coamenable. It is also
disconnected, since the Hopf inclusion of functions on $K$ gives $G\twoheadrightarrow K$, and the profinite group $K/K^0$
has a nontrivial finite quotient. We first give the reduction used in
both the cyclic and the higher-rank growth arguments. A group shift
over a finite group $F$ means a closed translation-invariant subgroup
of $F^D$, with pointwise multiplication.

\begin{lem}\label{shiftmodel}
Suppose $G$ is compact matrix and write $P=K/K^0$. The reverse action
of $K$ on $D$ is by a finite group of lattice automorphisms. Its kernel
$H$ contains $K^0$. The matched pair descends to $(D,P)$, and
$Q=H/K^0$ is an infinite open subgroup of $P$ on which $D$ acts by
automorphisms. As a $D$-system, $Q$ is isomorphic to a group shift over
a finite group.
\end{lem}
\begin{proof}
Let $S$ be the union of the $K$-orbits of a basis of $D$ and its
negatives. This generating set is finite. The matched-pair identity
applied to $(-s)+s$ gives $-\beta_k(s)=\beta_{\alpha_s(k)}(-s)$, so
$S$ is symmetric. For $s\in S$ one has
\[
 \beta_k(d+s)-\beta_k(d)=\beta_{\alpha_d(k)}(s)\in S
\]
Thus $\beta_k$ is an automorphism of the Cayley graph of $(D,S)$ fixing
zero. Here its linearity has an elementary proof. Each coordinate of the
identity map is harmonic since $S=-S$, and so is each coordinate of
$\beta_k$. Its difference along a basis vector is harmonic by translation
invariance and takes only finitely many values by the displayed identity.
At its maximum every neighbor has the same value. Connectedness makes
the difference constant, which proves linearity. This is also a case of
\cite{mmv}. There are finitely many possible maps because the basis is
sent into $S$. Substituting additivity into the identity for $\beta_k(d+s)$ gives
$\beta_{\alpha_d(k)}=\beta_k$.

The kernel $H$ of this right action is open, normal and $D$-invariant.
On $H$ the maps $\alpha_d$ are automorphisms, and
$\alpha_d(hk)=\alpha_d(h)\alpha_d(k)$ for $h\in H$. Since $K^0=H^0$,
this identity gives descent to $P$. The finite index of $H$ makes
$Q$ infinite and open in $P$.

Restriction to $H$ gives the ordinary crossed product
$C(H)\rtimes D$ as a quantum subgroup of $G$, hence it is compact matrix.
Write finitely many algebra generators of its coefficient algebra as
sums $\sum_j f_js_j$. All the functions so obtained lie in the
coefficient space of a finite-dimensional unitary representation $\rho$
of $H$. Take the algebra generated by the $D$-translates of those
coefficients and their adjoints. Crossing it with $D$ contains the
chosen generators, and taking the part of degree zero gives all of
$\mathcal R_H$. Hence
$h\mapsto(\rho(\alpha_d(h)))_{d\in D}$ embeds $H$ as a closed subgroup
of a compact Lie group shift.

By the descending chain theorem for compact Lie group shifts
\cite[Theorem~3.2]{ks}, the closed $D$-invariant subgroups of
$H$ satisfy the descending chain condition. It passes to $Q$ by
preimages. For $U$ open normal in the profinite group $Q$ set
$N_U=\bigcap_{d\in D}\alpha_d(U)$. The family of these invariant
subgroups is closed under finite intersections and its total intersection
is trivial. A minimal member is therefore trivial. For the corresponding
$U$, reading the quotient $Q\to Q/U$ at all translates embeds $Q$
as a closed group shift on the finite alphabet $Q/U$.
\end{proof}

We will use the following observation when locating a torus of a
quotient. For $E=nD$ and a closed $E$-invariant abelian subgroup
$A\subset Q$ we have
\[
 \mathcal R_G\ \supset\
 \mathcal R_P\rtimes_{alg} E
 \ \longrightarrow\
 \mathbb C[\widehat A\rtimes E]
\]
Indeed $E$ is characteristic in $D$ and the orbit indicators in
$\Delta(s_d)$ descend to $P$. Thus the middle algebra is the coefficient
Hopf $*$-algebra of $(E,P)$. The last arrow restricts functions to $A$,
is onto by Peter--Weyl, and respects the coproduct because $\beta$ is
trivial on $A$. Whenever $\widehat A\rtimes E$ is finitely generated,
the middle algebra can be replaced by a compact matrix coefficient
algebra. Lift group generators to it and take the finitely many coefficient
spaces of irreducibles needed to contain the lifts. Those spaces and
their adjoints generate a coefficient Hopf $*$-algebra of compact matrix
type still surjecting onto the group algebra. For $P=K$ and $E=D$ the
arrows already give a torus in $G$ itself.

We now specialize to $D=\mathbb Z$. Write $\mathbb Z_p$ for the cyclic group
of prime order $p$.

\begin{lem}\label{lampshift}
Let $X$ be an infinite closed shift-invariant subgroup of $F^{\mathbb Z}$,
where $F$ is a finite group and $T$ is the bilateral shift. For some
prime $p$, $X$ contains a closed invariant subgroup isomorphic, with its
shift, to $(\mathbb Z_p)^{\mathbb Z}$.
\end{lem}
\begin{proof}
Let $P_l$ be the group of words of length $l$ appearing in $X$. The
subgroups $\{a\in F:1^la\in P_{l+1}\}$ decrease with $l$ and stabilize,
say at $r\ge1$. The possible letters following a word of length $l\ge r$
form a coset of this stabilized group. They are contained in the coset
for the suffix of length $r$, so the two cosets coincide. Induction and
closedness now show that a configuration belongs to $X$ whenever all its
blocks of length $r+1$ occur in $X$.

As $X$ is infinite, $|P_l|$ is unbounded. For $l>2r$ with
$|P_l|>|F|^{2r}$, restriction to the first and last $r$ letters has
nontrivial kernel. Pad a nonidentity word in this kernel with identities
on both sides. Every block of length $r+1$ is admissible, so we have a
nonidentity element of finite support in $X$.

Among such elements choose one with the smallest support cardinality.
A suitable power $y$ has prime order $p$ and still has minimal support.
For $j\ne0$ the commutator $[y,T^jy]$ is supported in
$supp(y)\cap(supp(y)+j)$, strictly smaller
than $supp(y)$. It is therefore trivial. Thus all translates
of $y$ commute, even at sites where their supports overlap.

Arbitrary products of powers of these translates are defined coordinatewise,
as only finitely many factors contribute at each site. They give a continuous
shift-equivariant homomorphism $(\mathbb Z_p)^{\mathbb Z}\to X$ with nontrivial
compact image $A$. On taking Pontryagin duals, $\widehat A$ embeds as a
nonzero submodule of the regular module $R=\mathbb F_p[z,z^{-1}]$,
where $z$ acts by translation. Every nonzero ideal of this Laurent
polynomial ring is principal and is isomorphic to $R$ as an $R$-module.
Dualizing gives the required isomorphism with its shift.
\end{proof}

\begin{thm}\label{cycliclamps}
Let $G$ be a compact matrix bicrossed product of $\mathbb Z$ and a compact
group $K$ with infinitely many connected components. For some prime $p$
there is a compact matrix quantum quotient $L$ of $G$ such that
\[
 G\twoheadrightarrow L\supset\widehat\Lambda,\qquad \Lambda=\mathbb Z_p\wr\mathbb Z
\]
When $K$ is profinite one can take $L=G$. In all cases $\widehat G$
is amenable and has exponential growth.
\end{thm}

\begin{proof}
Use $P,Q$ from Lemma~\ref{shiftmodel}. Lemma~\ref{lampshift} gives a
closed shift-invariant subgroup $A\subset Q$ isomorphic with its
shift to $(\mathbb Z_p)^{\mathbb Z}$. Its character group consists of the finitely supported functions
from $\mathbb Z$ to $\mathbb Z_p$, with the dual shift. Taking $E=D$ and restricting functions on $P$ to $A$ gives the toral group
\[
 \widehat A\rtimes D=\left(\bigoplus_{n\in\mathbb Z}\mathbb Z_p\right)\rtimes\mathbb Z
                       =\mathbb Z_p\wr\mathbb Z
\]
One lamp and the shift generate this group, so the compact matrix
quotient obtained by lifting these generators suffices. For profinite $K$ we
have $K=P$, and the torus already lies in $G$.

Finally the shift can visit $n$ consecutive sites, at each site choosing
whether to switch the lamp on. This gives $2^n$ distinct group elements
of length at most $2n$. Lifts of the two generators have bounded length
in $\mathcal R_G$, so the same exponential lower bound holds
there after a linear change of radius. Amenability follows from that of
$\mathbb Z$ by \cite[Corollary~3.7]{fmp}.
\end{proof}

For the full shift $K=(\mathbb Z_p)^{\mathbb Z}$ with trivial reverse action, $G$
itself is the dual of $\mathbb Z_p\wr\mathbb Z$. Theorem~\ref{cycliclamps} shows
that every compact matrix bicrossed product in the stated class contains
this example at the level of tori of quotients.

\section{Lattice actions, growth and characters}\label{seclattice}

In two dimensions a group shift can be infinite while its only finitely
supported configuration is the identity. For instance, Ledrappier's
three-dot system \cite{ledrappier} consists of the arrays over $\mathbb F_2$
satisfying
\[
 x_{i,j+1}=x_{i,j}+x_{i+1,j}
\]
For a finitely supported configuration, take a zero row below its
support. The recurrence makes every row above it zero as well. Yet
projection onto one horizontal row gives the full binary shift. Indeed,
a predecessor of any two-sided row $a$ is obtained by fixing $b_0$ and
solving $b_i+b_{i+1}=a_i$ in both directions. We first explain what replaces the finite-support
argument for an arbitrary finite group alphabet.

\begin{lem}\label{abelianshift}
Let $D=\mathbb Z^r$ with $r\ge1$, let $F$ be a finite group and let $X$
be an infinite closed shift-invariant subgroup of $F^D$. Some sublattice
$nD$, $n\ge1$, preserves an infinite closed abelian subgroup of $X$.
\end{lem}
\begin{proof}
We induct on $|F|$. By passing to the image at one coordinate we can
suppose that all coordinate maps are onto $F$. If $F$ is abelian, take
$X$ itself. If $1<N<F$ is normal, put $Y=X\cap N^D$. When $Y$ is
infinite apply induction there. When $Y$ is finite the image of $X$ in
$(F/N)^D$ is infinite. By induction that image contains an infinite
abelian subgroup invariant under some $nD$ and closed in the image.
Let $B$ be its preimage in $X$. All commutators in $B$ lie in $Y$.
Choose an open normal subgroup $V$ of the profinite group $B$ with
$V\cap Y=1$. Then $[V,B]\subset V\cap Y=1$, so $Z(B)$ is open in $B$.
This center is infinite, closed, abelian and invariant under $nD$.

It remains to treat a nonabelian simple $F$. Put two sites in the same
class if their coordinate maps on $X$ have the same kernel. Within one
class the coordinates determine one another by automorphisms of $F$.
On distinct classes the values are independent. Indeed, on finitely many
sites Scott's lemma \cite[p.~328]{scott} writes the subdirect product as a
product of full diagonal subgroups. The diagonal blocks are exactly the
sites with equal kernels. Compactness then permits arbitrary choices on
one representative of each class. In particular a configuration can be set
equal to the identity off any union of classes while keeping its values
on that union.

Let $L\le D$ stabilize the class of $0$. The classes are indexed by
$D/L$, an infinite set since $X$ is infinite. Configurations supported
on the class of $0$ form a copy of $F$. The action of $L$ on this copy
has finite image in $Aut(F)$. Denote its kernel by $L_0$. Choose $n\ge1$ annihilating the torsion of the finitely
generated abelian group $D/L_0$. Then $n(D/L_0)$ is torsion-free,
and $L/L_0$ is finite. Hence $nD\cap L\subset L_0$.

Take a nonidentity element in this supported copy of $F$. Two of its
$nD$-translates either have disjoint supports or are equal, since a
translation in $nD$ stabilizing the class acts trivially there. There
are infinitely many supports because the image of $nD$ has finite
index in $D/L$. The closure of the group generated by these translates
is infinite, abelian and $nD$-invariant.
\end{proof}

\begin{lem}\label{latticelamps}
Let $E$ be a free abelian group of finite positive rank. Suppose $M$ is
an infinite discrete abelian group of bounded exponent, finitely generated
as a $\mathbb Z[E]$-module. Then $M\rtimes E$ contains $\mathbb Z_p\wr\mathbb Z$
for some prime $p$.
\end{lem}
\begin{proof}
If each basis element of $E$ had finite orbits on each of a finite set
of module generators, their full $E$-orbits would be finite by
commutativity. These finitely many elements would generate $M$ as a
group, making it finite. Thus some basis element $t$ has an infinite
orbit. Among elements with infinite $t$-orbit take $a$ of least order.
If its order is composite, take a prime $p$ dividing it. The orbit of
$pa$ is finite by minimality, so $b=(t^k-1)a$ is killed by $p$ for some
$k>0$. The orbit of $b$ is also finite by minimality. Its translates
generate a finite $t$-invariant subgroup, modulo which $t^k$ fixes $a$.
This would give $a$ a finite orbit. Hence $a$ has prime order $p$.

A nonzero Laurent polynomial relation over $\mathbb F_p$ between the
translates of $a$ would make their span a quotient of
$\mathbb F_p[z,z^{-1}]/(f)$ with $f\ne0$, a finite group. Thus the
translates are independent and $\langle a,t\rangle=\mathbb Z_p\wr\mathbb Z$.
\end{proof}

\begin{thm}\label{latticegrowth}
Let $G$ be a compact matrix bicrossed product of $D=\mathbb Z^r$, $r\ge1$,
and a compact group $K$ with infinitely many connected components.
There exist a compact matrix quantum quotient $L$ of $G$, a finitely
generated metabelian group $\Lambda$ and a prime $p$ such that
\[
 G\twoheadrightarrow L\supset\widehat\Lambda,\qquad
 \mathbb Z_p\wr\mathbb Z\ \subset\ \Lambda
\]
In particular $\widehat G$ is amenable and the tori of its untwisted
quotients detect its exponential growth.
\end{thm}
\begin{proof}

Let $P=K/K^0$ and take $Q\subset P$ from Lemma~\ref{shiftmodel}.
Identify $Q$ with a group shift in $F^D$ for a finite group $F$. Lemma~\ref{abelianshift} gives
a closed infinite abelian $A\subset Q$ invariant under $E=nD$.
Put $M=\widehat A$ and $\Lambda=M\rtimes E$. Since $E$ is characteristic
in $D$, restriction to $A$ gives the coefficient maps of Section~\ref{seccomponent},
\[
 \mathcal R_G\ \supset\ \mathcal R_P\rtimes_{alg}E
        \ \longrightarrow\ \mathbb C[\Lambda]
\]

The character group $M=\widehat A$ is infinite of exponent dividing the
exponent of $F$. Take one site in each coset of $E$ in $D$. At each of
these sites the image of $A$ is finite abelian. The characters of these
finitely many images and all their $E$-translates separate the points of
$A$, so they generate $M$ by Pontryagin duality. Thus $M$ is finitely
generated over $\mathbb Z[E]$, and $\Lambda=M\rtimes E$ is finitely
generated as a group. Lift a finite set of group generators to the middle algebra.
The finitely many coefficient spaces containing the lifts generate the required compact matrix quotient $L$ on adjoining their adjoints.

Lemma~\ref{latticelamps} gives the indicated subgroup of $\Lambda$.
Words in its lamp and shift already give exponential growth in $\Lambda$
and, by lifting generators, in $\mathcal R_G$. Coamenability
follows from amenability of $D$ by \cite[Corollary~3.7]{fmp}.
\end{proof}

For an ordinary crossed product there is also a shorter route through a
single lattice direction. It gives the lamplighter group itself as the
discrete group tested by the quotient.

\begin{cor}\label{ordinaryrow}
Under the assumptions of Theorem~\ref{latticegrowth}, suppose $D$ acts
on $K$ by automorphisms and the reverse action is trivial. There is a
compact matrix quantum quotient of $G$ containing the dual of $\mathbb Z_p\wr\mathbb Z$
for some prime $p$.
\end{cor}
\begin{proof}
Here $Q=P=K/K^0$. Embed $P$ in $F^D$ by Lemma~\ref{shiftmodel} and let $Y_i$ be its
image on the coordinate line $\mathbb Z e_i$. At least one $Y_i$ is
infinite. Otherwise the shift on each finite $Y_i$ has a common period
$m_i$ for its points. Translating the line shows that every configuration
in $P$ is periodic in each direction with period $m_i$, leaving only
finitely many choices in all.

For this infinite $Y_i$, restriction to the line is an equivariant group
surjection $P\to Y_i$ for $\mathbb Z e_i$. Pullback and the unitary $s_{e_i}$
give the Hopf $*$-subalgebra
$\mathcal R_{Y_i}\rtimes_{alg}\mathbb Z e_i\subset
\mathcal R_G$. Its compact quantum group is compact matrix,
since the functions at site zero and the shift generate the algebra.
Lemma~\ref{lampshift} supplies an invariant copy of $(\mathbb Z_p)^{\mathbb Z}$
in $Y_i$. Restricting to it gives the required torus.
\end{proof}

Theorem~\ref{latticegrowth} establishes Conjectures~\ref{twistamen}
and~\ref{twistgrowth} for the stated lattice bicrossed products.
Amenability is automatic by \cite[Corollary~3.7]{fmp}, and the
exponential-growth group already occurs with the trivial twist. The
inclusion $\mathbb Z_p\wr\mathbb Z\subset\Lambda$ here is an inclusion of
discrete groups. Corollary~\ref{ordinaryrow} obtains the lamplighter
dual itself as a torus in the ordinary crossed-product case.

We turn to characters. Pullback from the compact factor puts
$C(K)_{central}$ inside $C(G)_{central}$. If all tori kill a nonzero function in
this subalgebra, then joint character detection fails. For lattice
bicrossed products the converse also holds, even when the reverse
action is nontrivial.

\begin{thm}\label{bicrosscharacters}
Let $G$ be a compact matrix bicrossed product of $D=\mathbb Z^r$,
$r\ge1$, and a compact group $K$ with infinitely many connected
components. The tori of $G$ are jointly faithful on $C(G)_{central}$ if and only
if they are jointly faithful on $C(K)_{central}$.
\end{thm}
\begin{proof}
Suppose the tori are jointly faithful on $C(K)_{central}$, and fix a
nonzero $x\in C(G)_{central}$. We use right coefficients, writing an element
of the algebraic span of the characters as
\[
 y=\sum_{n\in D}s_n f_n,\qquad f_n\in\mathcal R_K
\]
For any character $\chi$, the sum of the diagonal coefficients gives
$\Delta(\chi)=\Delta^{op}(\chi)$, so the same equality holds for $y$.
Comparing the right coefficients of bidegree $(n,s)$ in the coproduct
formula of Section~\ref{secmatched}, we obtain, for $k,l\in K$,
\[
 1_{\{\beta_k(n)=s\}} f_n(kl)
   =1_{\{\beta_l(s)=n\}} f_s(lk)
\]
Take $l=1$ and $s=\beta_k(n)$. Since $\beta_1(s)=s$, a nonzero value
$f_n(k)$ forces $\beta_k(n)=n$. Taking $n=s=0$ gives $f_0(kl)=f_0(lk)$,
so $f_0$ is a class function. Both facts can also be read from the
character formula in \cite[Theorem~3.1]{fw}.

Let $E:C(K)\rtimes D\to C(K)$ take the coefficient at $0$. Since $D$
is amenable, the full and reduced crossed products agree and $E$ is
faithful. On the algebraic span just considered, $E$ takes values in
the class functions. By continuity the same holds on $C(G)_{central}$, which
is a $C^*$-algebra. In particular $E(x^*x)$ is a nonzero positive
class function. Choose a torus restriction
$\pi:C(G)\twoheadrightarrow C^*(\Gamma)$ on which it survives.
The image of $C(K)$ is $C(A)$ for a closed abelian subgroup $A\subset K$,
since restriction is a Hopf map into a cocommutative algebra.
Covariance makes $A$ invariant under each $\alpha_n$. We may therefore
choose $a\in A$ with $E(x^*x)(a)>0$.

By Lemma~\ref{shiftmodel}, $\sigma=\beta_a$ is a lattice automorphism
of finite order $m$, the image of $\beta$ is finite, and
$\beta_{\alpha_n(b)}=\beta_b$. Put
\[
 A_0=\{b\in A:\beta_b\in\langle\sigma\rangle\},\qquad
 N=1+\sigma+\cdots+\sigma^{m-1}\ \in End(D)
\]
Since the image of $\beta$ is finite, $A_0$ is open in $A$. It contains
$a$ and is invariant under
the lattice action. We have $N\beta_b=N$ for every $b\in A_0$, and
$Nn=mn$ on the subgroup fixed by $\sigma$. In particular $N$ is
injective on the lattice degrees that can have a nonzero coefficient
at $a$.

We first restrict the chosen torus over $A_0$. The projection
$p=1_{A_0}$ commutes with $C(A)$ and with every $\pi(s_n)$ by
invariance. These generate $C^*(\Gamma)$, so $p$ is central.
The Hopf inclusion $C(A)=C^*(\widehat A)\subset C^*(\Gamma)$ identifies
$\widehat A$ with a subgroup of $\Gamma$. Writing $F=A_0^\perp$ for the
finite subgroup of characters of $A$ trivial on $A_0$, we have
\[
 p=\frac1{|F|}\sum_{h\in F}h
\]
Centrality and linear independence of group elements give
$gFg^{-1}=F$ for every $g\in\Gamma$. The corner $pC^*(\Gamma)$ is
therefore the full group algebra of $\Gamma/F$. Indeed the unitary
representations on which $p$ acts as the identity are precisely those
trivial on $F$. Composing $\pi$ with the group quotient gives a torus
restriction $\pi_0$ with $\pi_0(C(K))=C(A_0)$, since compression on
$C(A)$ is restriction to $A_0$.

Write $v_n=\pi_0(s_n)$ and let $u_j$ denote the canonical unitaries of
$C^*(D)$. The following assignments give a map $\rho$ into
$C^*((\Gamma/F)\times D)$ by covariance.
\[
 f\longmapsto\pi_0(f)\otimes1,\qquad
 s_n\longmapsto v_n\otimes u_{Nn}
\]
It intertwines coproducts on $C(K)$. In the coproduct of $s_n$, the
summand with second degree $s$ has first-leg indicator
$1_{\{b:\beta_b(n)=s\}}$. After applying $\pi_0$ this indicator can
survive only if $s=\beta_b(n)$ for some $b\in A_0$. We then have
$Ns=Nn$, so the two added group factors in this summand are both
$u_{Nn}$. This proves compatibility with the coproduct on $s_n$ as well.
The image of the polynomial Hopf algebra is a Hopf $*$-subalgebra of
$\mathbb C[(\Gamma/F)\times D]$, hence is the group algebra of a
subgroup. Passing to its full completion gives a torus of $G$ through
which $\rho$ factors.

Finally extend evaluation at $a$ on $C(A_0)$ to a state of
$C^*(\Gamma/F)$. In its cyclic representation let $\xi$ be the unit
vector inducing this state. For $f\in C(K)$,
$\pi_0(f)\xi=f(a)\xi$, by applying the state to $|f-f(a)|^2$.
Use the regular representation on $\ell^2(D)$ for the other group
factor. Although $v_n$ need not be group-like, it is unitary. With the
right coefficients of $y$ we obtain
\[
 \rho(y)(\xi\otimes\delta_0)
       =\sum_n f_n(a)\,v_n\xi\otimes\delta_{Nn}
\]
The nonzero summands are orthogonal, since their degrees are fixed
by $\sigma$ and $N$ is injective there. We obtain
\[
 \|\rho(y)(\xi\otimes\delta_0)\|^2
       =\sum_n |f_n(a)|^2=E(y^*y)(a)
\]
The last equality follows by extracting degree zero in the product
using the right coefficients. Approximate $x$ in norm by such $y$.
Both sides are continuous, so the squared norm at $x$ equals
$E(x^*x)(a)>0$. The torus just constructed detects $x$. The reverse
implication follows from $C(K)_{central}\subset C(G)_{central}$.
\end{proof}

\section{Character detection and the quantum center}\label{seccenter}

The compact factor in Theorem~\ref{bicrosscharacters} belongs to the
given bicrossed-product presentation. To obtain a character test defined
for $G$ alone, consider its quantum center $Z(G)$. We recall the
description by the chain group \cite{chirvasitu,kss}. Let $\Gamma$ be
the group with one generator $g_u$ for each $u\in Irr(G)$ and relations
\[
 g_w=g_u g_v \qquad\text{whenever } w\in u\otimes v
\]
In particular $g_1=e$ and $g_{\bar u}=g_u^{-1}$. The group $\Gamma$
grades the polynomial Hopf algebra by assigning degree $g_u$ to every
matrix coefficient $u_{ij}$. Restriction to the quantum center is
given on matrix coefficients by $u_{ij}\mapsto\delta_{ij}g_u$, so
$Z(G)=\widehat\Gamma$ \cite[Remark~2.7]{kss}. An irreducible
representation descends to $G/Z(G)$ precisely when its degree is $e$.
Taking the norm closure of the pulled-back characters in $C(G)$ gives
\[
 C(G/Z(G))_{central}=\overline{span}\{\chi_u:u\in Irr(G),\ g_u=e\}
\]
Throughout this section restrictions on this subalgebra are taken to
the tori of $G$.

Write $I$ for the common kernel of the torus restrictions on $C(G)_{central}$.
Define $P_e$ on the span of the characters by keeping only the terms
of degree $e$. To see its extension to the norm closure, let $U_g$ be
the canonical unitary in $C^*(\Gamma)$ and define the grading map by
\[
 \delta(u_{ij})=u_{ij}\otimes U_{g_u},\qquad
 P_e=(id\otimes\tau)\delta
\]
Here $\tau(U_g)$ is one for $g=e$ and zero otherwise. The map $\delta$
is the coproduct followed by restriction of the second factor to the
center, with values in the spatial tensor product
$C(G)\otimes C^*(\Gamma)$. Slicing by the trivial representation
$U_g\mapsto1$ gives the identity, so $\delta$ is injective. The
slice defining $P_e$ is positive and contractive. On $C(G)_{central}$ it is a
conditional expectation onto $C(G/Z(G))_{central}$.

\begin{prop}\label{neutralprojection}
For any compact matrix quantum group $G$ of Kac type, $P_e(I)$ is
contained in $I$. If the chain group is amenable, every nonzero
$x\in I$ gives a nonzero element $P_e(x^*x)$ in $I\cap C(G/Z(G))_{central}$.
The same conclusion holds for any nonzero finite linear combination
of irreducible characters in $I$, for an arbitrary chain group.
\end{prop}
\begin{proof}
Fix a torus restriction $\pi:C(G)\twoheadrightarrow C^*(\Lambda)$.
For $a$ homogeneous of degree $g$, the assignment
\[
 a\longmapsto \pi(a)\otimes U_g\ \in\ \mathbb C[\Lambda\times\Gamma]
\]
on the polynomial algebra preserves multiplication and the involution.
It also preserves coproducts, since both legs of the coproduct of a
matrix coefficient of an irreducible of degree $g$ have degree $g$.
The image is a Hopf $*$-subalgebra of a group algebra, hence the group
algebra of a subgroup. The resulting map to its full completion is a
torus restriction of $G$. Its composition with the natural map to
$C^*(\Lambda)\otimes C^*(\Gamma)$ agrees with
$(\pi\otimes id)\delta$ on the polynomial algebra and
therefore on $C(G)$. For $x\in I$ the enlarged torus kills $x$.
Slicing gives $\pi(P_e x)=0$. This holds for every $\pi$.

When $\Gamma$ is amenable, the regular trace $\tau$ is faithful on
$C^*(\Gamma)$, so $id\otimes\tau$ is faithful on positive
elements of the spatial tensor product. Injectivity of $\delta$ now
gives $P_e(x^*x)\ne0$ for $x\ne0$. For a finite sum $x=\sum_g x_g$
of homogeneous characters or their linear combinations we have directly
\[
 P_e(x^*x)=\sum_g x_g^*x_g
\]
This is nonzero as soon as some $x_g$ is nonzero. In both cases $x^*x$
belongs to $I$, and the first assertion finishes the proof.
\end{proof}

\begin{cor}\label{classicalneutral}
Let $G\twoheadrightarrow K$ be a quantum quotient with $K$ a classical compact group. Whenever the
tori of $G$ kill a nonzero $f\in C(K)_{central}$, they also kill a nonzero
element of $C(G/Z(G))_{central}$. In particular New conjecture~\ref{newchar} holds
for both crossed-product classes in Theorem~\ref{introchar}.
\end{cor}
\begin{proof}
Regard representations of $K$ as representations of $G$ by pullback.
For $u,v\in Irr(K)$ the tensor products $u\otimes v$ and $v\otimes u$
are isomorphic. Any common irreducible constituent therefore gives
$g_u g_v=g_v g_u$ in the chain group of $G$. The subgroup $A$ generated
by the degrees coming from $K$ is abelian. Consequently
\[
 \delta(C(K)_{central})\ \subset\ C(K)_{central}\otimes C^*(A)
\]
Here $C^*(A)$ embeds in $C^*(\Gamma)$, and the restriction of $\tau$
to it is the faithful regular trace of the abelian group $A$.
Thus $P_e(f^*f)$ is nonzero. Proposition~\ref{neutralprojection} shows
that it belongs to $I\cap C(G/Z(G))_{central}$.

For an ordinary crossed product with amenable discrete factor,
Theorem~\ref{classcriterion} and its proof give a nonzero invisible
class function on $K$ whenever joint detection on $C(G)_{central}$ fails. For a
lattice bicrossed product with infinitely many compact-factor components,
Theorem~\ref{bicrosscharacters} gives precisely that implication. Applying
the first assertion proves the character conjecture in both cases.
\end{proof}

A counterexample to New conjecture~\ref{newchar} would therefore have
a nonamenable chain group. In such an example every nonzero character polynomial would survive on
some torus, but a nonzero norm limit of character polynomials would
vanish on all of them. When $Z(G)$ is trivial, $G/Z(G)=G$ and the
conjecture holds immediately.

\section{Examples and partition calculations}\label{sectests}

An old torus is still available for the two analytic tests. For
characters the verification concerns the norm closure in $C(G)_{central}$.

\begin{prop}\label{transfer}
Within the compact matrix Kac setting, each verification of Old
conjecture~\ref{oldamen} or~\ref{oldgrowth} gives a verification of
New conjecture~\ref{newamen} or~\ref{newgrowth}, respectively.
Joint detection of $C(G)_{central}$ implies New conjecture~\ref{newchar}.
This last conjecture also holds whenever $\widehat G$ is amenable or
$G$ has an irreducible fundamental representation.
\end{prop}
\begin{proof}
Every old torus occurs in the new family by taking the identity
quotient and the trivial twist. Thus a witness to failure of either
analytic property still works, while Proposition~\ref{twistnecessary}
proves the forward implications. Joint detection of $C(G)_{central}$ is the
conclusion of New conjecture~\ref{newchar}, with or without connectedness.

For the two other character assertions, let $\Gamma$ be the chain group.
Its dual is the central torus $Z(G)\subset G$, so amenability of
$\widehat G$ gives amenability of $\Gamma$ by
Proposition~\ref{twistnecessary}. If a fundamental representation $u$
is irreducible, every irreducible occurs in a tensor word in $u,\bar u$.
Its degree is thus a power of $g_u$, and $\Gamma$ is cyclic.
In both cases apply Proposition~\ref{neutralprojection}.
\end{proof}

In particular, all three new conjectures hold for the classical compact
matrix groups and group duals of \cite[Proposition~4.1]{bp}. On a
classical group each element belongs to the closure of the cyclic
subgroup it generates, so restrictions even separate all continuous
functions. A group dual is its own torus.

All three also hold for every compact quantum subgroup of $O_N^*$.
Indeed, polynomial growth of the dual of $O_N^*$ \cite[Section~9]{halforth}
passes to the dual of each such subgroup \cite[Proposition~4.1]{dpr}
and implies amenability \cite[Theorem~4.6]{dpr}.
Propositions~\ref{twistnecessary} and~\ref{transfer} now apply.
This covers the half-classical examples of \cite[Theorem~5.5]{bp},
including $[H]$ from the crossed-product construction of \cite{halfhopf},
and works for the new tests without a connectedness assumption.

\topic{Orthogonal easy quantum groups}
For New conjecture~\ref{newchar}, even the partitions describing the
higher tensor powers are unnecessary in the following case.

\begin{thm}\label{easycharacters}
If $S_N\subset G\subset O_N^+$ in the permutation representation,
then New conjecture~\ref{newchar} holds. In particular it holds for every
orthogonal easy quantum group.
\end{thm}
\begin{proof}
For $N\ge2$, the commutant of $S_N$ on $\mathbb C^N$ is spanned by the
identity and the projection onto the constant vectors. Indeed a commuting
matrix has one common diagonal entry and one common off-diagonal entry.
The commutant of the fundamental $u$ of $G$ is a unital $*$-subalgebra
of this two-dimensional algebra. If it consists of scalars, $u$ is
irreducible and Proposition~\ref{transfer} applies.

Otherwise $u$ has two irreducible summands on the constants and their
orthogonal complement. Both summands are self-conjugate, since their
projections are real and $u$ is orthogonal. Their degrees generate the
chain group and square to the identity. This group is therefore a
quotient of the amenable group $\mathbb Z_2*\mathbb Z_2$. Apply
Proposition~\ref{neutralprojection}. The case $N=1$ is immediate.
For an orthogonal easy quantum group the inclusion of $S_N$ follows
because every partition map intertwines the permutation action
\cite{liberation}.
\end{proof}

We can now account for the four free families of \cite[Theorem~4.4]{bp}.
The preceding theorem gives New conjecture~\ref{newchar} for $O_N^+$,
$S_N^+$ and $H_N^+$, and the fundamental representation of $U_N^+$
is irreducible already on $U_N$. For $N\ge3$ the standard diagonal
groups of $O_N^+,H_N^+$ are $\mathbb Z_2^{*N}$, while that of $U_N^+$
is $F_N$. For $S_N^+$ with $N\ge5$, a pair and a triple of points
give $\mathbb Z_2*\mathbb Z_3$. These groups witness nonamenability
and exponential growth in the indicated ranks, as does $F_2$ for
$U_2^+$. All remaining ranks have polynomial growth and amenable dual
\cite[Theorem~4.4]{bp}. Thus the two analytic verifications transfer
through Proposition~\ref{transfer}. For the character verification
above, no passage from polynomial injectivity to norm injectivity is
needed. The same transfer applies wherever the diagonal groups of the
other orthogonal easy families in \cite[Section~6]{bp} are known to be
nonamenable. Theorem~\ref{easycharacters} supplies the character
assertion for these families as well.

\topic{Torus relations from partitions}
For an orthogonal easy quantum group and a real orthogonal basis
change $Q$, take a partition map $T_p$ between tensor powers of
lengths $k,l$. Set
\[
 T_p^Q=Q^{\otimes l}T_p(Q^t)^{\otimes k}
\]
On the diagonal torus the intertwiner equation reads
\[
 (T_p^Q)_{j,i}\,g_{i_1}\cdots g_{i_k}
 =g_{j_1}\cdots g_{j_l}\,(T_p^Q)_{j,i}
\]
Whenever the scalar coefficient is nonzero, we can cancel it and identify
the two words. It suffices to impose these relations on any family
generating the tensor category
\cite[Section~6]{bp}. For unitary categories one also uses the conjugate
basis on conjugate tensor factors and inverse letters there. For
New conjectures~\ref{newamen} and~\ref{newgrowth}, a compact matrix
quantum quotient selects a full tensor subcategory closed under
subobjects, conjugates and finite sums. A dimension-preserving unitary
fiber functor realizes the allowed twist \cite{bdv}. One may then use the
same coefficient rule on the resulting concrete intertwiners. A single
presentation exhibiting a group without the required property suffices.

Here is a small example where the basis matters. Write $D_\infty=
\mathbb Z_2*\mathbb Z_2$ and take
\[
 Q=\begin{pmatrix}a&-b\\ b&a\end{pmatrix},\qquad a^2+b^2=1
\]
\begin{prop}\label{ranktwotori}
For $H_2^+$ the three possibilities for $\Gamma_Q$ are
\[
 \Gamma_Q\simeq
 \begin{cases}
 D_\infty & ab=0\cr
 \mathbb Z_2\times\mathbb Z_2 & |a|=|b|\cr
 \mathbb Z_2 & ab\ne0,\ |a|\ne|b|
 \end{cases}
\]
The same list covers all real orthogonal bases up to row signs.
\end{prop}
\begin{proof}
The pair partition makes the two generators $g,h$ involutions.
The four-block fixes $e_1^{\otimes4}+e_2^{\otimes4}$.
These two tensors present the orthogonal cubic relations
\cite{liberation,reflection}. Bending two legs of the four-block gives the projection onto the span of $e_i\otimes e_i$. Intertwining
this projection with $u^{\otimes2}$ says exactly that distinct entries
in any row or column have product zero.

If $ab=0$ the only supported four-letter words are $g^4,h^4$.
Otherwise the coefficient of $ghgh$ is $2a^2b^2$, so the two
involutions commute. A word with three copies of one letter and one
of the other has coefficient $\pm ab(a^2-b^2)$. When this is zero
the remaining relations have even multiplicity of each letter, and
commuting involutions satisfy them all. When it is nonzero they force
$g=h$, and one involution satisfies every relation. Row signs do
not change the support of either tensor.
\end{proof}

The original basis gives $D_\infty$. The normalized Fourier matrix gives
\[
 Q=\frac{1}{\sqrt2}\begin{pmatrix}1&1\\1&-1\end{pmatrix},\qquad
 \Gamma_Q\simeq\mathbb Z_2\times\mathbb Z_2
\]
In the Fourier basis we can check commutativity directly. Put $w=gh$.
Since $g$ and $h$ are involutions, we have
\[
 u_{11}=\frac{g+h}{2},\qquad u_{11}^2=\frac{2+w+w^{-1}}4
\]
The cubic relations make $u_{11}^2$ a projection. Squaring the second
expression and equating it to itself gives $w^2+w^{-2}=2$.
Linear independence in the group algebra forces $w^2=1$, so $g$ and $h$
commute.
This corrects the Fourier-basis calculation in the arXiv version of
\cite[proof of Theorem~4.4]{bp}.

Computing the two tensor supports over
the integers, \texttt{tori\_partitions.py} finds 4, 4 and 632 primitive
signed pairs in the respective cases for $\max(|a|,|b|)\le16$
before normalization.

\begin{prop}\label{reflectiontest}
New conjectures~\ref{newchar}--\ref{newgrowth} hold for the free
complex reflection groups $H_N^{s+}$, for $N\ge2$ and $s\ge3$.
Their standard diagonal group is the free product of $N$ copies of
$\mathbb Z_s$, and their chain group is $\mathbb Z_s$.
\end{prop}
\begin{proof}
In the defining presentation of \cite{reflection}, the entries of $u$
are normal partial isometries satisfying $u_{ij}^{s}=u_{ij}u_{ij}^*$.
Setting all off-diagonal entries to zero leaves arbitrary unitaries
$g_i$ with $g_i^s=1$. Conversely any such family gives a diagonal
solution of all the relations. In the partition description a block of
$s$ equal colors imposes $g_i^s=1$ in this basis. The resulting group contains $\mathbb Z_s*\mathbb Z_s$. The kernel of
the map from their free product onto $\mathbb Z_s^2$ acts freely on the
free product tree. The quotient graph is the complete bipartite graph on $s+s$ vertices.
Counting its edges and vertices gives the free rank of the kernel
\[
 s^2-2s+1=(s-1)^2\ge2
\] This one
torus therefore verifies both analytic tests.

The classical diagonal phase matrices and the permutation matrices act
irreducibly together, so $u$ is irreducible. Its degree generates the
chain group. The tensor $\sum_i e_i^{\otimes s}$ is fixed by the block
of size $s$, so $g_u^s=1$. Scalar matrices of order $s$ give a central
subgroup on which $u$ has the defining degree of order $s$. This proves
the claimed equality, and Proposition~\ref{transfer} proves the
character assertion.
\end{proof}

The script counts words by syllable length, using the generators and their
inverses. For $\mathbb Z_3*\mathbb Z_3$ it finds the following ball volumes
at radii $0,\ldots,6$
\[
 1,\ 5,\ 13,\ 29,\ 61,\ 125,\ 253
\]
At the same radii in $D_\infty$ the volumes are $1,3,5,7,9,11,13$.

\topic{A test with nonamenable chain group}
The following dilation argument also gives the full norm in the
single-factor calculation of \cite[Theorem~4.4]{bp}.

\begin{thm}\label{freeunitarytest}
Let $C(G)$ be the unital full free product of the coefficient
algebras of $U_{N_1}^+,\ldots,U_{N_r}^+$ with their usual coproducts,
where $r\ge1$ and $N_i\ge2$. The diagonal torus of the block
fundamental representation is isometric on $C(G)_{central}$. All three new
conjectures hold, and the chain group is $F_r$.
\end{thm}
\begin{proof}
Write $\chi_i$ for the fundamental character of the $i$-th factor.
The fusion rules for a single factor \cite[Theorem~6.20]{qg} and the
alternating-word description of the free product irreducibles
\cite[Theorem~3.10]{wangfree} give $C(G)_{central}=C^*(\chi_1,\ldots,\chi_r)$.
Let $g_{ij}$ be the $j$th diagonal generator of the $i$th factor. The
restriction of its character is
\[
 \pi(\chi_i)=g_{i1}+\cdots+g_{iN_i}
\]
All these generators are free, both within each factor and across
different factors.

Fix a representation of $C(G)$ on a Hilbert space $E$. The operators
$T_i=\chi_i/N_i$ in this representation are contractions. For each
$i$ form the unitary on $E\oplus E$
\[
 V_i=\begin{pmatrix}
 T_i&\sqrt{1-T_iT_i^*}\\
 \sqrt{1-T_i^*T_i}&-T_i^*
 \end{pmatrix}
\]
Summing the conjugates over the $N_i$-th roots of unity cancels the
off-diagonal entries, giving
\[
 D_z=\begin{pmatrix}1&0\\0&z\end{pmatrix},\qquad
 \sum_{z^{N_i}=1}D_zV_iD_z^*
 =N_i\begin{pmatrix}T_i&0\\0&-T_i^*\end{pmatrix}
\]
Assign these unitaries $D_zV_iD_z^*$ to
the generators of the $i$-th block. Freeness gives a representation of
the full toral group algebra on $E\oplus E$. The first summand reduces
all the character images and recovers the original $\chi_i$
simultaneously. Taking the supremum over representations proves equality
of norms for every $*$-polynomial in the characters and hence on $C(G)_{central}$.

The toral group $F_{\sum_i N_i}$ witnesses nonamenability and exponential
growth. The $r$ fundamental degrees generate the chain group and map to
the free generators $z_i$ of $F_r$ under $u_i\mapsto z_i I_{N_i}$.
This quotient is central, as one checks on the coefficients of $u_i$.
Thus the chain group is $F_r$.
\end{proof}


\topic{Twisting the free product}
The preceding isometry uses a particular torus. To follow that torus
under a twist we keep track of the bases of the fundamental blocks.

\begin{cor}\label{freetwist}
For the free product $G$ of Theorem~\ref{freeunitarytest}, every unitary
dual $2$-cocycle is a coboundary. In every twist, the diagonal torus of
the corresponding fundamental blocks has group $F_{N_1+\cdots+N_r}$,
and restriction to this torus is isometric on the full character algebra.
\end{cor}
\begin{proof}
The cocycle defines a dimension-preserving unitary fiber functor on
$Rep(G)$. On $U_{N_i}^+$ such a functor is classified
by an invertible matrix $M$ of size $N_i$, up to left and right
unitaries \cite[Theorem~6.2 and Remark~6.5]{bdv}. For $P=M^*M$ the two
trace conditions give
\[
 Tr(P)=Tr(P^{-1})=N_i,\qquad Tr(P+P^{-1}-2I)=0
\]
Each eigenvalue $\lambda>0$ of $P$ contributes
$\lambda+\lambda^{-1}-2=(\lambda-1)^2/\lambda$ to the last trace.
All contributions are nonnegative, so every eigenvalue equals one and $P=I$.
It follows that $M$ is equivalent to the identity and the restricted
functor is unitarily monoidally isomorphic to the forgetful functor.

The representation category of the free product has the universal
property of the free product of the factor categories
\cite[Proposition~5.1]{valvekens}. Hence the factor isomorphisms
extend to a unitary monoidal isomorphism on all of $Rep(G)$.
On a tensor word its value is the tensor product of the factor
isomorphisms, composed with the tensor structure maps of the functors.
By \cite[Proposition~4.5]{bdv} this says precisely that the cocycle
is a coboundary.

The induced Hopf $*$-isomorphism matches each fundamental block up to
unitary change of basis. Any such changes can be absorbed independently
by automorphisms of the original free product, since
$u_i\mapsto W_i u_i W_i^*$ preserves biunitarity and has an inverse.
The isomorphism then identifies the corresponding diagonal quotients
and preserves character norms by Theorem~\ref{freeunitarytest}.
\end{proof}

For an explicit twist take two $U_2^+$ factors, with fundamental
matrices $u,v$ and toral letters $g_1,g_2,h_1,h_2$, respectively.
Define $p:F_4\to A=\mathbb Z_2^2$ on the generators by
\[
 p(g_1)=(1,0),\qquad p(h_1)=(0,1),\qquad p(g_2)=p(h_2)=0
\]
On $A$ take the bicharacter $\sigma(x,y)=(-1)^{x_1y_2}$. Its values on
the two basis elements depend on their order. Since every coboundary on
$A$ is symmetric, $\sigma$ represents a nontrivial class. Composing
with $p$ gives a cocycle on $F_4$. The diagonal restriction
$\mathcal R_G\to\mathbb C[F_4]$ then gives the cocycle on
the dual of $G$.

For coefficients $a,b$ of respective left and right $A$-weights
$(x,x')$ and $(y,y')$, the twisted product is
\[
 a\star b=\sigma(x,y)\overline{\sigma(x',y')}\,ab
\]
For example, the order of the two coefficients affects the sign
\[
 u_{12}\star v_{11}=-u_{12}v_{11},\qquad
 v_{11}\star u_{12}=v_{11}u_{12}
\]
For a word in the free generators and their inverses put
\[
 b(\ell_1\cdots\ell_k)=
 (-1)^{\displaystyle\sum_{i<j}p(\ell_i)_1p(\ell_j)_2}
\]
An inverse pair contributes zero internally, since a letter has weight
on at most one coordinate. Its terms with every other letter cancel.
Thus $b$ is defined on $F_4$, and splitting the sum between words gives
\[
 b(st)=b(s)b(t)\sigma(p(s),p(t))
\]
For a coefficient of left and right $F_4$-degrees $(s,t)$, multiplication
by $b(s)\overline{b(t)}$ untwists the product. This follows by
applying the preceding identity to the left and right degrees separately. It fixes the
entries of $u$ and $v$, since $b$ is one on single letters.
In $\mathbb C[F_4]$, a group element has equal left and right degrees.
Hence diagonal restriction kills coefficients with $s\ne t$. For
$s=t$ the untwisting sign is one. Thus in this example the original restriction
map itself survives on the common polynomial coalgebra, with
\[
 \chi_u\longmapsto g_1+g_2,\qquad \chi_v\longmapsto h_1+h_2
\]
The script \texttt{tori\_twists.py} checks this formula on all 69,905
formal coefficient words through length four, including adjoints and
the empty word. All 341 character $*$-words through that length give
identical toral expansions. Corollary~\ref{freetwist} gives the isometry
on the norm closure.

\phantomsection\addcontentsline{toc}{section}{References}
\begin{thebibliography}{99}
\small\setlength{\itemsep}{0pt}\setlength{\parskip}{0pt}
\bibitem{qg} T. Banica, Introduction to quantum groups, Springer (2023).
\bibitem{generic} T. Banica, Symmetries of a generic coaction,
\emph{Math. Ann.} \textbf{314} (1999), 763--780.
\bibitem{bbd} T. Banica, J. Bhowmick and K. De Commer, Quantum isometries and group dual subgroups, \emph{Ann. Math. Blaise Pascal} \textbf{19} (2012), 1--27.
\bibitem{bp} T. Banica and I. Patri, Maximal torus theory for compact quantum groups,
\emph{Illinois J. Math.} \textbf{61} (2017), 151--170.
\bibitem{liberation} T. Banica and R. Speicher, Liberation of orthogonal Lie groups, \emph{Adv. Math.} \textbf{222} (2009), 1461--1501.
\bibitem{reflection} T. Banica and R. Vergnioux, Fusion rules for quantum reflection groups, \emph{J. Noncommut. Geom.} \textbf{3} (2009), 327--359.
\bibitem{halforth} T. Banica and R. Vergnioux, Invariants of the half-liberated orthogonal group, \emph{Ann. Inst. Fourier} \textbf{60} (2010), 2137--2164.
\bibitem{bdv} J. Bichon, A. De Rijdt and S. Vaes, Ergodic coactions with large multiplicity and monoidal equivalence of quantum groups, \emph{Comm. Math. Phys.} \textbf{262} (2006), 703--728.
\bibitem{halfhopf} J. Bichon and M. Dubois-Violette, Half-commutative orthogonal Hopf algebras, \emph{Pacific J. Math.} \textbf{263} (2013), 13--28.
\bibitem{behy} M. Brannan, F. Elzinga, S.J. Harris and M. Yamashita, Crossed product equivalence of quantum automorphism groups of finite dimensional $C^*$-algebras, \emph{Int. Math. Res. Not.} \textbf{2023} (2023), 17749--17787.
\bibitem{btd} T. Br\"ocker and T. tom Dieck, Representations of compact Lie groups, Graduate Texts in Mathematics \textbf{98}, Springer (1985).
\bibitem{chirvasitu} A. Chirvasitu, Centers, cocenters and simple quantum groups,
\emph{J. Pure Appl. Algebra} \textbf{218} (2014), 1418--1430.
\bibitem{dpr} A. D'Andrea, C. Pinzari and S. Rossi, Polynomial growth of discrete quantum groups, topological dimension of the dual and $*$-regularity of the Fourier algebra, \emph{Ann. Inst. Fourier} \textbf{67} (2017), 2003--2027.
\bibitem{fmp} P. Fima, K. Mukherjee and I. Patri, On compact bicrossed products,
\emph{J. Noncommut. Geom.} \textbf{11} (2017), 1521--1591.
\bibitem{fw} P. Fima and H. Wang, Rapid decay and polynomial growth for bicrossed products, \emph{J. Noncommut. Geom.} \textbf{15} (2021), 1105--1128.
\bibitem{goswami} D. Goswami, A note on outer quantum automorphisms of finite dimensional von Neumann algebras, \emph{Indian J. Pure Appl. Math.} \textbf{55} (2024), 922--930.
\bibitem{kss} P. Kasprzak, A. Skalski and P.M. So{\l}tan, The canonical central exact sequence for locally compact quantum groups,
\emph{Math. Nachr.} \textbf{290} (2017), 1303--1316.
\bibitem{ks} B. Kitchens and K. Schmidt, Automorphisms of compact groups, \emph{Ergodic Theory Dynam. Systems} \textbf{9} (1989), 691--735.
\bibitem{kyed} D. Kyed, $L^2$-Betti numbers of coamenable quantum groups, \emph{M\"unster J. Math.} \textbf{1} (2008), 143--180.
\bibitem{ledrappier} F. Ledrappier, Un champ markovien peut \^etre d'entropie nulle et m\'elangeant, \emph{C. R. Acad. Sci. Paris S\'er. A} \textbf{287} (1978), 561--563.
\bibitem{mmv} D.W. Morris, J. Morris and G. Verret, Isomorphisms of Cayley graphs on nilpotent groups, \emph{New York J. Math.} \textbf{22} (2016), 453--467.
\bibitem{mrozinski} C. Mrozinski, Quantum automorphism groups and $SO(3)$-deformations, \emph{J. Pure Appl. Algebra} \textbf{219} (2015), 1--32.
\bibitem{scott} L.L. Scott, Representations in characteristic $p$, \emph{Proc. Sympos. Pure Math.} \textbf{37} (1980), 319--331.
\bibitem{valvekens} M. Valvekens, Some remarks on free products of rigid $C^*$-2-categories, arXiv:2008.12994 (2020).
\bibitem{wangfree} S. Wang, Free products of compact quantum groups, \emph{Comm. Math. Phys.} \textbf{167} (1995), 671--692.
\bibitem{wang} S. Wang, Tensor products and crossed products of compact quantum groups,
\emph{Proc. London Math. Soc.} \textbf{71} (1995), 695--720.
\bibitem{simple} S. Wang, Simple compact quantum groups I, \emph{J. Funct. Anal.} \textbf{256} (2009), 3313--3341.
\bibitem{wcmp} S.L. Woronowicz, Compact matrix pseudogroups, \emph{Comm. Math. Phys.} \textbf{111} (1987), 613--665.
\bibitem{wtk} S.L. Woronowicz, Tannaka--Krein duality for compact matrix pseudogroups. Twisted $SU(N)$ groups, \emph{Invent. Math.} \textbf{93} (1988), 35--76.
\end{thebibliography}
\end{document}
